\documentclass[landscape, 8pt, a4paper, oneside]{extarticle}
\AtBeginDocument{\fontsize{8}{8}\selectfont}

\usepackage{teamnote}
\usepackage{multicol}
\usepackage{listings}
\usepackage{pdflscape}
\usepackage{iftex}


\ifPDFTeX
  \DeclareUnicodeCharacter{2060}{\nolinebreak}
  \DeclareUnicodeCharacter{2212}{\nolinebreak}
  \usepackage[utf8]{inputenc}
  \usepackage[T5]{fontenc}
  \usepackage[vietnamese,english]{babel}
\else
  \usepackage{fontspec}
  \usepackage[vietnamese,english]{babel}
  \setmainfont{DejaVu Serif}
  \setsansfont{DejaVu Sans}
  \setmonofont{DejaVu Sans Mono}
\fi
\usepackage{graphicx}
\usepackage{tikz}
\usepackage{geometry}
\geometry{landscape, a4paper, left=0.25in, right=0.25in, top=0.35in, bottom=0.25in, headheight=14pt}

\ShowUsage
\ShowComplexity
\HideAuthor

\begin{document}
\selectlanguage{vietnamese}

\begin{titlepage}
    \begin{tikzpicture}[remember picture,overlay]
        \node[opacity=0.20,inner sep=0pt] at (current page.center)
        {\includegraphics[height=0.7\paperheight]{logo_uit.png}};
    \end{tikzpicture}

    \centering
    {\fontsize{16}{19} \bfseries ICPC Notebook -- \teamname \par}
    \vspace{0.15cm}
    {\fontsize{10}{12} \bfseries \schoolname \par}
    {\fontsize{9}{11} \authorname \quad | \quad \today\par}
    \vspace{0.2cm}
    \begin{multicols}{3}
        \raggedright
        \tableofcontents
    \end{multicols}
\end{titlepage}

\begin{multicols}{3}

\section*{Tối ưu hằng số}
\addcontentsline{toc}{section}{Tối ưu hằng số}
\begin{lstlisting}[language=C++]
#include <bits/stdc++.h>
#pragma GCC optimize("O3,unroll-loops")
// Chi bat target neu CPU judge ho tro.
// #pragma GCC target("avx2,bmi,bmi2,lzcnt,popcnt")
\end{lstlisting}
Đặt sau include, trước các hàm. Chỉ bật \texttt{target} khi CPU judge hỗ trợ; đo lại vì \texttt{O3}/unroll có thể chậm hơn \texttt{O2}. \texttt{Ofast} nới quy tắc số thực.

\textbf{Đo}: xác định nghẽn tính toán, bộ nhớ hay I/O; cùng input/compiler, thay một yếu tố mỗi lần, đo vài lượt. Test lại với brute force; máy local và judge có thể khác tốc độ.

\textbf{Cache/thứ tự duyệt}: ưu tiên mảng liên tiếp; chỉ số cuối ở vòng trong. Nếu đổi chiều lưu trữ, đổi đồng bộ mọi truy cập; chỉ đổi thứ tự xử lý khi không vi phạm phụ thuộc dữ liệu.

\textbf{Kiểu dữ liệu}: dùng \texttt{int32\_t} cho chỉ số/giá trị vừa 32-bit khi template có macro \texttt{int}; tổng/tích lớn dùng 64-bit hoặc \texttt{\_\_int128}, ép kiểu trước phép nhân. Không dựa vào tràn số có dấu.

\textbf{Cấp phát}: tái sử dụng buffer; \texttt{reserve(n)} trước nhiều lần \texttt{push\_back}. Reserve chỉ đặt capacity; cần \texttt{resize} trước khi truy cập phần tử bằng chỉ số. Nhiều node nhỏ: cân nhắc pool/mảng node thay cấp phát từng node.

\textbf{Array vs vector}: mảng C/\texttt{std::array<T,N>} có kích thước cố định lúc biên dịch; \texttt{vector<T>} có kích thước động. Cả ba lưu liên tiếp (trừ \texttt{vector<bool>}); truy cập \texttt{[]} đều $O(1)$, thường tương đương khi bật tối ưu. Vector tốn thêm cấp phát/quản lý capacity; \texttt{vector<int>(n)} khởi tạo $n$ số $0$, mảng số nguyên cục bộ không tự khởi tạo. Biết cận cố định, tái dùng nhiều: mảng; cần đúng kích thước input/thay đổi size: vector. Mảng lớn nên đặt global/static để tránh tràn stack; vector nhiều chiều lồng nhau không bảo đảm toàn khối liên tiếp, có thể dùng vector phẳng.

\textbf{Template vector nhiều chiều}: \texttt{Vec<D,T>}, $D$ là số chiều, $T$ là kiểu phần tử. Constructor nhận kích thước từng chiều, cuối cùng là giá trị khởi tạo (có thể bỏ, mặc định \texttt{T()}). Đây vẫn là vector lồng nhau, không phải một khối phẳng.
\inputminted[]{cpp}{source/template/multidimensional.cpp}
\begin{lstlisting}[language=C++]
Vec<2, int> mat(rows, cols, 0);
Vec<3, int> cube(rows, cols, depth, -1);
cube[xx][yy][zz] = 5;
\end{lstlisting}

\textbf{Container}: ID nhỏ/nén được thì dùng mảng/vector; cần thứ tự mới dùng map/set. Với \texttt{unordered\_map}, thử \texttt{reserve} và chỉnh load factor; không mặc định hash luôn nhanh hơn. Đồ thị lớn: cân nhắc mảng cạnh liên tiếp.

\textbf{Hash table}: GNU \texttt{gp\_hash\_table} là lựa chọn đáng đo thử khi nghẽn hash; không suy rộng tốc độ từ benchmark của người khác. Collision có thể làm tổng thời gian thành $O(n^2)$; \texttt{reserve} không ngăn hack. Cân nhắc hash trộn tốt với seed ngẫu nhiên; chỉ XOR/cộng một hằng ngẫu nhiên chưa đủ.

\textbf{Tránh sao chép/gọi gián tiếp}: truyền container bằng \texttt{const\&}; duyệt phần tử lớn bằng \texttt{const auto\&}. Hàm gọi rất nhiều lần: ưu tiên hàm thường/template/lambda cụ thể nếu không cần \texttt{std::function}.

\textbf{Vòng lặp nóng}: tính biểu thức bất biến/base một lần, tái dùng kết quả, tiền xử lý đại lượng dùng nhiều lần. Bỏ cập nhật đóng góp $0$ nếu không có tác dụng phụ; dừng sớm khi đã đủ kết luận. Không tự unroll/branchless nếu chưa đo.

\textbf{Số học}: tránh \texttt{pow} cho bình phương số nguyên; dùng nhân với kiểu đủ rộng. Với $0\le a,b<mod$, cộng modulo bằng cộng rồi trừ $mod$ nếu cần; tổng trung gian phải không tràn. Giảm số lần chia/modulo trong vòng nóng nếu có công thức tương đương đúng.

\textbf{Bit song song}: giao hai tập bằng \texttt{(maskA \& maskB).count()}; subset sum boolean với trọng số không âm: \texttt{can[0]=1}, mỗi phần tử \texttt{can |= can << weight}. Cấp đủ số bit cho miền tổng cần xét; mỗi phép trên bitset vẫn quét nhiều word, không phải $O(1)$.

\textbf{Duyệt bit 1}: dùng \texttt{\_\_builtin\_popcountll} cho mask 64-bit. Mask không âm: \texttt{while(mask)\{ bit=\_\_builtin\_ctzll(mask); mask\&=mask-1; \}}; không gọi \texttt{ctz/clz} với $0$. GNU bitset có \texttt{\_Find\_first()}, \texttt{\_Find\_next(pos)}; hết bit 1 trả \texttt{size()}. Đây là extension, vẫn có thể quét các word trống.

\textbf{Reset}: chỉ xóa vị trí đã chạm hoặc dùng timestamp nếu mỗi test chỉ dùng ít phần tử trên miền lớn; xử lý giá trị mặc định và tràn timestamp. Có thể tái sử dụng vùng nhớ giữa các pha.

\textbf{I/O}: \texttt{ios::sync\_with\_stdio(false); cin.tie(nullptr);}, dùng \texttt{'\textbackslash n'} thay \texttt{endl}. Dữ liệu lớn: thử \texttt{fread}/\texttt{fwrite}, xuất theo khối; tránh trộn C/C++ I/O sau khi tắt sync. Interactive vẫn phải flush khi gửi câu hỏi.




\section{Cấu trúc dữ liệu}
\noindent\textbf{Quy ước C++17}: Chỉ số từ 0, đoạn $[l,r)$. File nguồn có sẵn header chuẩn và namespace; phần bọc lặp lại được ẩn khi in.

\Algorithm{Iterative Segment Tree}
{Segment tree trên monoid có thứ tự: gán điểm, gộp đoạn, tìm biên; bản dual cộng đoạn/truy vấn điểm}
{$\mathcal{O}(\log n)$ mỗi truy vấn/cập nhật, $\mathcal{O}(n)$ dựng}
{compactcpp}{source/cpp/data_structure/iterative_segment_tree.cpp}


\Algorithm{Lazy Segment Tree (lưu ý cài đặt)}
{Quy tắc lazy propagation, hợp thành tag (Affine, Range Set) và Tree Walk $O(\log n)$}
{}
{}{}

\begin{itemize}
    \item \textbf{Quy tắc push/pull}: Gọi \texttt{push(id, l, r)} đầu \texttt{update}/\texttt{query}; gọi \texttt{pull(id)} (\texttt{st[id] = merge(st[2*id], st[2*id+1])}) sau khi cập nhật các con.
    \item \textbf{Nhiều lazy tag}:
    \begin{itemize}
        \item \textbf{Affine ($x \mapsto a \cdot x + b$)}: Áp dụng tag mới $(c,d)$ sau $(a,b) \implies c(ax+b)+d = (ca)x + (cb+d)$. \\ Tag: \texttt{laz\_a = laz\_a * c; laz\_b = laz\_b * c + d;} \\ Node: \texttt{val = val * c + (r - l + 1) * d;}
        \item \textbf{Range Set và Range Add}: Range Set ($x \leftarrow v$) ghi đè Range Add trước đó $\implies$ đặt lại \texttt{add = 0}.
    \end{itemize}
    \item \textbf{Tree Walk ($O(\log n)$)}: Nếu con trái thỏa ($\max \ge X$ hoặc tổng $\ge S$), đẩy lazy và đi trái; nếu không, điều chỉnh mục tiêu rồi đi phải.
    \item \textbf{Giới hạn}: Cấp phát $4N$ khi đánh số từ 1 (\texttt{id << 1} và \texttt{id << 1 | 1}).
\end{itemize}



\Algorithm{Sparse Segment Tree}
{Segment tree động, chỉ cấp phát node khi cần}
{$\mathcal{O}(\log n)$ mỗi truy vấn}
{compactcpp}{source/cpp/data_structure/sparse_segment_tree.cpp}

\Algorithm{Cartesian Tree}
{Cây nhị phân có thứ tự heap theo giá trị và thứ tự BST theo chỉ số}
{$\mathcal{O}(n)$ dựng bằng monotonic stack}
{}{}

\begin{itemize}
    \item \textbf{Tính chất}: Min-heap theo giá trị ($val[\text{parent}] \le val[\text{child}]$) + duyệt in-order cho chỉ số mảng gốc ($1..n$). Gốc là giá trị nhỏ nhất. Cây con $u$ ứng với đoạn cực đại $[L_u, R_u]$ mà $A[u]$ là nhỏ nhất ($size(u) = R_u - L_u + 1$).
    \item \textbf{$\mathcal{O}(n)$ Dựng bằng stack}: Giữ nhánh phải trong stack. Với $i = 1..n$: pop $u$ có $A[u] > A[i]$; node pop cuối thành $left[i] = u$; đỉnh còn lại có $right[\text{top}] = i$. Push $i$.
    \item \textbf{Ứng dụng chính}:
    \begin{itemize}
        \item \textbf{RMQ $\iff$ LCA}: Giá trị nhỏ nhất trên đoạn $\min_{l \le k \le r} A[k] = A[\text{LCA}(l, r)]$.
        \item \textbf{Hình chữ nhật lớn nhất trong histogram}: $\max_u (A[u] \cdot size(u))$.
        \item \textbf{D\&C $\to$ Tree DP}: Chuyển D\&C trên mảng, chia tại phần tử nhỏ nhất, thành DP trên cây con.
    \end{itemize}
\end{itemize}

\Algorithm{Persistent Segment Tree}
{Persistent segment tree: cập nhật điểm, truy vấn các phiên bản cũ}
{$\mathcal{O}(\log n)$ mỗi truy vấn/cập nhật}
{compactcpp}{source/cpp/data_structure/persistent_segment_tree.cpp}

\Algorithm{Parallel Binary Search}
{Binary search song song offline cho nhiều truy vấn đơn điệu}
{$\mathcal{O}((N + Q) \log (\text{range}) \log N)$}
{compactcpp}{source/cpp/data_structure/parallel_binary_search.cpp}

\Algorithm{Wavelet Tree}
{Wavelet tree: phần tử thứ $k$, tần suất và đếm giá trị trong $[L, R]$ trên đoạn}
{$\mathcal{O}(\log (\text{alphabet}))$ mỗi truy vấn}
{compactcpp}{source/cpp/data_structure/wavelet_tree.cpp}

\Algorithm{Segment Tree Beat}
{Segment Tree Beats: chmin/chmax, cộng và truy vấn tổng trên đoạn}
{$\mathcal{O}((n+q)\log^2 n)$ khấu hao}
{compactcpp}{source/cpp/data_structure/segment_tree_beats.cpp}

\Algorithm{Kd Tree}
{Cây $k$ chiều tĩnh: điểm gần nhất và truy vấn hộp đóng song song trục}
{$O(n)$ mỗi truy vấn trong trường hợp xấu nhất; cắt tỉa không gian trên cây cân bằng}
{compactcpp}{source/cpp/data_structure/kd_tree.cpp}

\Algorithm{BIT2D}
{Fenwick 2D: tổng hình chữ nhật và cập nhật điểm}
{$\mathcal{O}(\log^2 n)$ mỗi truy vấn}
{compactcpp}{source/cpp/data_structure/fenwick_2d.cpp}

\Algorithm{Treap}
{BST ngẫu nhiên: split/merge mảng và thao tác theo khóa ẩn}
{$\mathcal{O}(\log n)$ kỳ vọng mỗi thao tác}
{compactcpp}{source/cpp/data_structure/implicit_treap.cpp}

\Algorithm{Splay Tree / Link - cut Tree}
{Link-Cut Tree: truy vấn đường đi động, link, cut và đổi gốc trên rừng}
{$\mathcal{O}(\log n)$ khấu hao mỗi thao tác}
{compactcpp}{source/cpp/data_structure/link_cut_tree.cpp}

\Algorithm{Farach - Colton - Bender}
{$\langle \mathcal{O}(n), \mathcal{O}(1) \rangle$ LCA qua Euler Tour, RMQ $\pm 1$ và kỹ thuật Four Russians}
{$\mathcal{O}(n)$ dựng, $\mathcal{O}(1)$ truy vấn}
{compactcpp}{source/cpp/data_structure/farach_colton_bender.cpp}

\Algorithm{Trie}
{Trie nhị phân multiset 64-bit: thêm, xóa, đếm, XOR lớn nhất}
{$\mathcal{O}(\text{length})$ mỗi truy vấn}
{compactcpp}{source/cpp/data_structure/binary_trie.cpp}

\Algorithm{Lichao Tree}
{Li Chao Tree: thêm đường thẳng online, truy vấn min/max}
{$\mathcal{O}(\log n)$ mỗi truy vấn/thêm}
{compactcpp}{source/cpp/data_structure/li_chao.cpp}

\Algorithm{Line Container}
{Convex hull của các đường thẳng động, hệ số góc và điểm truy vấn $x$ tùy ý}
{$\mathcal{O}(\log n)$ khấu hao mỗi thao tác}
{compactcpp}{source/cpp/data_structure/line_container.cpp}

\Algorithm{Arpa Trick}
{Truy vấn min/max đoạn offline bằng DSU}
{$\mathcal{O}(n \alpha(n))$}
{compactcpp}{source/cpp/data_structure/arpa_trick.cpp}

\Algorithm{Chia căn (Square Root Decomposition)}
{Chia căn: mảng chia block và thuật toán Mo theo thứ tự Hilbert}
{$\mathcal{O}(n \sqrt{n})$}
{compactcpp}{source/cpp/data_structure/block_array_mo.cpp}

\Algorithm{Hash Map (SplitMix64)}
{Hash map khóa không dấu, hash SplitMix64 với seed do người gọi chọn}
{$\mathcal{O}(1)$ kỳ vọng mỗi lần tra cứu}
{compactcpp}{source/cpp/data_structure/hash_map.cpp}

\Algorithm{Erasable Priority Queue}
{Priority queue max/min, xóa lazy phần tử tùy ý}
{$\mathcal{O}(\log n)$ mỗi thao tác}
{compactcpp}{source/cpp/data_structure/erasable_priority_queue.cpp}

\Algorithm{Cấu trúc mở rộng (Binary Doubling)}
{Kỹ thuật gộp block theo lũy thừa 2: thêm online cho cấu trúc tĩnh (Aho, CHT)}
{$\mathcal{O}(A \log n)$ thêm khấu hao, với cấu trúc tĩnh dựng được trong $O(A)$}
{compactcpp}{source/cpp/data_structure/expandable.cpp}

\Algorithm{Interval Set}
{Phân hoạch đoạn: gán đoạn, gộp các đoạn kề nhau cùng giá trị.}
{$O((k+1)\log s)$ mỗi lần gán, $k$ đoạn bị xóa, $s$ đoạn đang lưu}
{compactcpp}{source/cpp/data_structure/interval_set.cpp}


\section{Đồ thị}

% Reference: https://math.mit.edu/~goemans/18453S17/matroid-intersect-notes.pdf
\subsection{Matroid: nhận diện và các loại}
Ví dụ: chọn cạnh không tạo cycle là \textbf{graphic matroid}; thêm quota mỗi màu là giao với \textbf{partition matroid}. Một matroid dùng greedy; hai matroid cần trao đổi cả chuỗi phần tử.

\textbf{Định nghĩa}: $M=(E,\mathcal I)$, $\mathcal I$ là các tập independent:
\begin{itemize}
    \item $\varnothing\in\mathcal I$; $A\in\mathcal I, B\subseteq A\Rightarrow B\in\mathcal I$.
    \item \textbf{Exchange}: $A,B\in\mathcal I, |A|<|B|\Rightarrow\exists x\in B\setminus A: A+x\in\mathcal I$.
\end{itemize}
\textbf{Rank} $r(X)=\max_{S\subseteq X,S\in\mathcal I}|S|$; \textbf{basis}: independent tối đại, luôn có $r(E)$ phần tử. \textbf{Circuit}: tập dependent tối tiểu. Greedy xét trọng số giảm dần, thêm nếu còn independent, cho maximum-weight basis; nếu không buộc đủ rank thì bỏ trọng số âm.

\begin{itemize}
    \item \textbf{Uniform}: $|S|\le k$.
    \item \textbf{Partition}: $E$ chia nhóm rời nhau $E_i$, $|S\cap E_i|\le c_i$. Oracle: đếm quota.
    \item \textbf{Graphic}: $E$ là cạnh, $S$ là forest. Oracle: DSU; circuit là cycle.
    \item \textbf{Linear}: $E$ là vector trên cùng field, $S$ độc lập tuyến tính. Oracle: Gauss / XOR basis trên $\mathbb F_2$.
    \item \textbf{Laminar}: quota $|S\cap A|\le c_A$, các $A$ đôi một lồng nhau hoặc rời nhau. Ví dụ quota trên mọi prefix.
    \item \textbf{Transversal}: $E$ là phía trái của bipartite graph; $S$ ghép được vào các đỉnh phải khác nhau. Oracle: matching phủ $S$.
\end{itemize}
\textbf{Cue}: xóa phần tử vẫn hợp lệ chưa đủ; phải có exchange. Các \emph{tập cạnh matching} không là matroid: trên path ba cạnh, matching gồm cạnh giữa không thêm được cạnh nào từ matching gồm hai cạnh ngoài.

\subsection{Matroid intersection (maximum cardinality)}
Hai matroid trên \textbf{cùng $E$}; cần oracle $I_j(S)$ kiểm tra independent. Bắt đầu $S=\varnothing$. Mỗi vòng dựng exchange graph với $x\notin S$, $y\in S$:
\[
\begin{array}{c|l}
\text{Cạnh} & \text{Điều kiện} \\
s\to x & I_1(S+x) \\
y\to x & I_1(S-y+x) \\
x\to y & I_2(S-y+x) \\
x\to t & I_2(S+x)
\end{array}
\]
\textbf{BFS} lấy shortest path $s\leadsto t$: đảo trạng thái chọn của mọi phần tử trên path (bỏ $s,t$). Khi đó $|S|$ tăng $1$ và vẫn independent ở cả hai matroid. Dựng lại graph rồi lặp; không còn path $\Rightarrow$ tối ưu.

\textbf{Principle khi code}: giữ nguyên $S$ khi gọi oracle/dựng graph; chỉ cập nhật sau khi tìm \emph{shortest} path. Không dùng một path tùy ý. Với $N=|E|$, $r=\min(r_1(E),r_2(E))$, oracle tốn $T$: cách dựng trực tiếp dùng $O(Nr^2)$ oracle calls, thời gian $O(Nr^2(T+1))$. Weighted intersection cần thuật toán riêng; thay BFS bằng heap chưa đủ.

\textbf{Mô hình}: rainbow forest = graphic + quota màu; bipartite matching = hai partition theo hai đầu cạnh. Giao hai matroid thường không còn là matroid nên greedy trực tiếp không đúng.

\textbf{Interface}: \texttt{MatroidIntersection(n, oracle1, oracle2).solve()} trả về các ID được chọn. Oracle nhận \texttt{const vector<int>\&} gồm ID của tập cần kiểm tra, trả về \texttt{bool}; không giữ tham chiếu vào tập này.
\inputminted{compactcpp}{source/cpp/graph/matroid_intersection.cpp}


\subsection{Đếm thứ tự topo trên cây/rừng}
Cạnh hướng cha $\to$ con: số thứ tự topo là $\displaystyle \frac{n!}{\prod_v sz(v)}$, với $sz(v)$ là kích thước subtree kể cả $v$. Ghép các thứ tự con bằng multinomial; đảo toàn bộ cạnh vẫn cùng số cách. Không áp dụng cho DAG tổng quát.


\Algorithm{Minimum Cost Flow}
{Min-cost $k$-flow: chuyển $k$ đơn vị qua đồ thị có capacity và chi phí}
{$O(nm + A m\log n)$, $A$ lần tăng luồng}
{compactcpp}{source/cpp/graph/min_cost_flow.cpp}

\Algorithm{MinCut - MaxFlow}
{Dinic: số đường đi/đường đi không chung đỉnh lớn nhất từ source đến sink.
\begin{itemize}
    \item \textbf{Tô hai màu lưới}: Chia ô theo $(r + c) \bmod 2 \implies$ đồ thị hai phía. Dùng lát domino, đặt quân không tấn công nhau, bắn tia laser trên lưới.
    \item \textbf{Tách đỉnh (Vertex Splitting) ($u_{in} \to u_{out}$)}: Giới hạn capacity đỉnh hoặc buộc đi qua đỉnh (capacity $c_u$, chi phí $w_u$).
    \item \textbf{Max-Weight Closure (Project Selection)}:
    \begin{itemize}
        \item Quy tắc: chọn $u$ thì phải chọn $v$ (cạnh có hướng $u \to v$ với $\text{cap} = \infty$).
        \item Trọng số dương $W_u > 0$: thêm cạnh $S \to u$ với $\text{cap} = W_u$.
        \item Trọng số âm $W_v < 0$: thêm cạnh $v \to T$ với $\text{cap} = -W_v$.
        \item $\text{Max Profit} = \sum_{W_u > 0} W_u - \text{MinCut}(S, T)$.
    \end{itemize}
    \item \textbf{Phân nhóm với hai lựa chọn}: Gán $i$ vào $S$ (chi phí $A_i$) hoặc $T$ ($B_i$). Nếu $i \in S, j \in T$, phạt $C_{i,j}$. Thêm $S \to i$ ($B_i$), $i \to T$ ($A_i$), $i \to j$ ($C_{i,j}$). Tổng chi phí $=$ min-cut.
    \item \textbf{Circulation có cận dưới/trên $[L_e, R_e]$}: Thay cạnh $e=(u,v)$ bằng capacity $R_e-L_e$. Đặt $d(u)=\sum_{in}L-\sum_{out}L$. Thêm super-source $S^*$, super-sink $T^*$. Với $d(u)>0$, thêm $S^*\to u$ capacity $d(u)$; với $d(u)<0$, thêm $u\to T^*$ capacity $-d(u)$. Khả thi $\iff$ mọi cạnh từ $S^*$ bão hòa.
\end{itemize}
}
{$\mathcal{O}(EV^2)$}
{compactcpp}{source/cpp/graph/max_flow_dinic.cpp}

\Algorithm{Bipartite Matching}
{Hopcroft-Karp: vertex cover, path cover không chung đỉnh/tổng quát, antichain (Dilworth).
\begin{itemize}
    \item \textbf{Định lý K\H{o}nig}: $\text{Max Matching} = \text{Min Vertex Cover}$.
    \item \textbf{Max Independent Set}: $|V| - \text{Min Vertex Cover} = |V| - \text{Max Matching}$.
    \item \textbf{Min Edge Cover} (không có đỉnh cô lập): $|V| - \text{Max Matching}$.
    \item \textbf{Path cover nhỏ nhất trên DAG}:
    \begin{itemize}
        \item \textit{Không chung đỉnh}: Tách $u \to u_{out}, u_{in}$; dựng matching hai phía $u_{out} \to v_{in}$ cho cạnh $u \to v$. Kết quả: $|V| - \text{Max Matching}$.
        \item \textit{Tổng quát (được chung đỉnh)}: Tính transitive closure (Floyd-Warshall/reachability), rồi giải path cover không chung đỉnh.
    \end{itemize}
    \item \textbf{Định lý Dilworth}: Trong poset hữu hạn (reachability trên DAG), kích thước \textbf{maximum antichain} (tập lớn nhất gồm các đỉnh không đến được nhau) bằng \textbf{minimum path cover}.
    \item \textbf{Erd\H{o}s--Gallai}: $n-1\ge d_1\ge\dots\ge d_n\ge0$ là dãy bậc đồ thị đơn vô hướng iff $\sum d_i$ chẵn và $\sum_{i=1}^{k}d_i\le k(k-1)+\sum_{i=k+1}^{n}\min(d_i,k)$ ($\forall 1\le k\le n$). Dựng hình $\to$ \textbf{Havel--Hakimi}: lấy $\max d$, nối $d$ đỉnh lớn kế tiếp rồi trừ 1.
    \item \textbf{Gale--Ryser}: Hai dãy nguyên không âm $a_1\ge\cdots\ge a_n$, $b_1,\dots,b_m$ là bậc hai phía của đồ thị đơn iff $\sum a_i=\sum b_j$ và $\sum_{i=1}^{k}a_i\le\sum_{j=1}^{m}\min(b_j,k)$ với mọi $1\le k\le n$. Dùng cho ma trận $0/1$ có tổng hàng/cột cho trước.
    \item \textbf{Định lý Hall}: Đồ thị hai phía $(X,Y)$ có matching phủ hết $X \iff \forall S \subseteq X, |N(S)| \ge |S|$.
    \item \textbf{Kích thước maximum matching qua Hall}: $|X| - \max_{S \subseteq X} (|S| - |N(S)|)$.
\end{itemize}
}
{$\mathcal{O}(E\sqrt{V})$}
{compactcpp}{source/cpp/graph/hopcroft_karp.cpp}

\Algorithm{Hungarian Matching}
{Ghép cặp hoàn hảo chi phí nhỏ nhất trên đồ thị hai phía}
{$\mathcal{O}(V^3)$}
{compactcpp}{source/cpp/graph/hungarian.cpp}

\Algorithm{General Matching}
{Edmonds Blossom: maximum matching trên đồ thị tổng quát.
\begin{itemize}
    \item \textbf{Tutte Matrix \& Rabin-Vazirani ($O(V^3)$)}: Với $(u,v)\in E$ ($u<v$), đặt $T_{u, v} = r \in [1, p-1]$, $T_{v, u} = -r \pmod p$. Kích thước $= \frac{1}{2}\text{rank}(T)$ (Gauss mod $p$).
    \begin{itemize}
        \item \textit{Khôi phục Rabin-Vazirani}: Tính $A=T^{-1}\pmod p$ (trên cơ sở $2K\times2K$). Cạnh $(u,v)$ thuộc một matching $\iff A_{u,v}\ne0$. Chọn cặp đó, ghép lại rồi cập nhật $A$ trong $O(V^2)$ bằng Schur complement: $A'_{x, y} = A_{x, y} + \frac{A_{x, u}A_{y, v} - A_{x, v}A_{y, u}}{A_{u, v}} \pmod p$. Lặp cho $K$ cặp $\implies$ tổng $O(V^3)$.
    \end{itemize}
    \item \textbf{Quy về max $b$-matching} ($\deg(i) \le a_i$): Đỉnh $i \to a_i$ slot ($a_i \leftarrow \min(a_i, \deg(i))$). Cạnh $e=(u,v) \to L_e, R_e$, nối $L_e - R_e$, $L_e - \text{slot}(u)$, $R_e - \text{slot}(v)$. Bỏ $e$: match $L_e-R_e$ ($+1$). Chọn $e$: match $L_e, R_e$ với slot ($+2$, lợi thêm $+1$). Đồ thị phụ $|V'| \le 4M \implies$ Blossom.
\end{itemize}
}
{$\mathcal{O}(V^3)$}
{compactcpp}{source/cpp/graph/general_matching.cpp}

\Algorithm{Biconnected Component}
{BCC: khớp, cầu và thành phần song liên thông}
{$\mathcal{O}(V + E)$}
{compactcpp}{source/cpp/graph/bcc.cpp}

\Algorithm{Functional Cycle}
{Floyd Tortoise and Hare: tìm chu trình trong functional graph $f(x)$}
{$O(\text{dist} + \text{len})$}
{compactcpp}{source/cpp/graph/functional_cycle.cpp}

\Algorithm{Strongly Connected Components}
{Tarjan SCC trên đồ thị có hướng}
{$\mathcal{O}(V + E)$}
{compactcpp}{source/cpp/graph/scc_tarjan.cpp}

\Algorithm{Xóa offline (DSU Rollback)}
{DSU rollback kết hợp D\&C trên segment tree}
{$\mathcal{O}(T(n) \log n)$}
{compactcpp}{source/cpp/graph/dsu_rollback.cpp}

\Algorithm{2-SAT}
{Giải 2-SAT bằng đồ thị kéo theo (implication graph) và SCC.
\begin{itemize}
    \item \textbf{Ít nhất một biến đúng ($X \lor Y$)}: Thêm các cạnh $\neg X \to Y$ và $\neg Y \to X$.
    \item \textbf{Nếu $X$ thì $Y$ ($X \implies Y$)}: Thêm các cạnh $X \to Y$ và $\neg Y \to \neg X$.
    \item \textbf{Đúng một biến đúng ($X \text{ XOR } Y = 1$)}: $(X \lor Y) \land (\neg X \lor \neg Y)$. Các cạnh: $X \leftrightarrow \neg Y$ và $\neg X \leftrightarrow Y$.
    \item \textbf{$X$ và $Y$ phải bằng nhau ($X \text{ XOR } Y = 0$)}: $(X \lor \neg Y) \land (\neg X \lor Y)$. Các cạnh: $X \leftrightarrow Y$ và $\neg X \leftrightarrow \neg Y$.
    \item \textbf{Buộc $X$ đúng/sai}: Thêm cạnh $\neg X \to X$ (buộc đúng), hoặc $X \to \neg X$ (buộc sai).
    \item \textbf{Nhiều nhất một biến trong $\{X_1, X_2, \dots, X_k\}$ đúng}: Thêm biến prefix $P_i = P_{i-1} \lor X_i$. Các cạnh: $X_i \to P_i$, $P_{i-1} \to P_i$, và $X_i \to \neg P_{i-1}$ (nếu $X_i$ đúng thì mọi biến trước phải sai).
\end{itemize}
}
{$\mathcal{O}(n + m)$}
{compactcpp}{source/cpp/graph/two_sat.cpp}

\Algorithm{Floyd}
{Đường đi ngắn nhất mọi cặp/transitive closure/đường đi min-max}
{$O(n^3)$}
{compactcpp}{source/cpp/graph/floyd_warshall.cpp}

\Algorithm{Eulerian Path}
{Hierholzer: tìm đường đi/chu trình Euler.
\begin{itemize}
    \item \textbf{Nối từ/chuỗi (De Bruijn/domino)}: Nối từ sao cho đuôi từ trước khớp đầu từ sau. Đỉnh là \textit{trạng thái} (prefix/suffix dài $L-1$ hoặc nửa domino); cạnh có hướng là \textit{đối tượng}. Đường đi Euler đi qua mỗi đối tượng đúng một lần.
    \item \textbf{Định hướng cạnh vô hướng}: Định hướng cạnh để cân bằng bậc vào/ra (ví dụ $|in-out|\le1$). Thêm cạnh giả (hoặc nối đỉnh bậc lẻ tới đỉnh giả) để mọi bậc chẵn; tìm chu trình Euler rồi định hướng theo lượt duyệt.
    \item \textbf{Phân rã chu trình không chung cạnh}: Phân rã thành chu trình không chung cạnh. Khả thi $\iff$ đồ thị liên thông và mọi đỉnh bậc chẵn. Tìm chu trình Euler rồi tách các chu trình.
    \item \textbf{Duyệt lưới bằng đồ thị hàng-cột}: Bài ma trận di chuyển luân phiên hàng/cột: dựng đồ thị hai phía gồm hàng và cột. Ô $(r,c)$ là cạnh nối hàng $r$ với cột $c$. Đường đi Euler duyệt các ô theo hàng-cột-hàng.
\end{itemize}
}
{$\mathcal{O}(m)$}
{compactcpp}{source/cpp/graph/euler_path.cpp}

\Algorithm{Định lý Vizing}
{Tô cạnh đồ thị đơn bằng tối đa $\Delta+1$ màu}
{$O(nm)$ thời gian, $O(n^2)$ bộ nhớ}
{compactcpp}{source/cpp/graph/vizing.cpp}

\Algorithm{Kiểm tra hai phía online}
{Kiểm tra hai phía online khi thêm cạnh bằng parity DSU}
{$\mathcal{O}(\alpha(n))$}
{compactcpp}{source/cpp/graph/online_bipartite.cpp}

\Algorithm{Gomory-Hu Tree}
{Cây cut tương đương cho max-flow/min-cut mọi cặp}
{$O(n \cdot \text{MaxFlow})$}
{compactcpp}{source/cpp/graph/gomory_hu.cpp}

\Algorithm{Critical (Dominator Tree)}
{Lengauer-Tarjan: immediate dominator từ nguồn $s$}
{$O(m \log n)$}
{compactcpp}{source/cpp/graph/dominator_tree.cpp}

\Algorithm{Prufer}
{Mã hóa/giải mã cây có nhãn gồm $n$ đỉnh}
{$O(n \log n)$}
{compactcpp}{source/cpp/graph/prufer.cpp}

\Algorithm{Steiner Tree}
{Steiner tree nhỏ nhất nối $k$ đỉnh terminal}
{$O(n \cdot 3^k + m \cdot 2^k \log m)$}
{compactcpp}{source/cpp/graph/steiner_tree.cpp}

\Algorithm{Stable Marriage}
{Gale-Shapley giải bài ghép đôi ổn định (Stable Marriage)}
{$O(n^2)$}
{compactcpp}{source/cpp/graph/stable_marriage.cpp}


\section{Tree}

\Algorithm{Tree Isomorphism}
{Tree hashing: Check isomorphic rooted trees or unrooted trees (using centroids as roots)}
{$\mathcal{O}(n \log n)$}
{}{source/pseudocode/tree_isomorphism.txt}

\Algorithm{Centroid Decomposition}
{Decompose tree into $O(\log n)$ depth hierarchy for divide-and-conquer path queries}
{$\mathcal{O}(n \log n)$}
{}{source/pseudocode/centroid_decomposition.txt}

\Algorithm{Heavy-Light Decomposition}
{Decompose tree into heavy chains for $O(\log^2 n)$ path queries/updates and $O(\log n)$ subtree queries}
{$\mathcal{O}(n \log^2 n)$}
{}{source/pseudocode/hld.txt}



\section{String}

\Algorithm{String Problem Solving Strategies}
{String problem taxonomy, periodicity, LCP, distinct substrings, palindromes \& key theorems}
{}
{}{}

\begin{itemize}
    \item \textbf{Problem Taxonomy \& Core Approches}:
    \begin{itemize}
        \item \textbf{Longest Common Prefix with Whole String / Single Pattern Match}: Use \textbf{Z-Algorithm} or \textbf{KMP}.
        \item \textbf{Multiple Dictionary Matching / DP on String Sets}: Use \textbf{Aho-Corasick Automaton}.
        \item \textbf{Distinct Substrings}: Suffix Array ($\frac{n(n+1)}{2} - \sum LCP_i$) or SAM ($\sum len[u] - len[link[u]]$).
        \item \textbf{$k$-th Substring / LCS of Multiple Strings}: Use \textbf{Suffix Automaton (SAM)} or \textbf{Suffix Array + LCP}.
        \item \textbf{Palindromic Substrings}: Use \textbf{Manacher} (radius centered at every index) or \textbf{EERTREE / PAM} (palindromic tree).
        \item \textbf{Lexicographically Minimal String Rotation}: Run \textbf{Lyndon Factorization (Duval)} on $S + S$.
        \item \textbf{Fast Substring Equivalence $O(1)$ / 2D Strings}: Use \textbf{Rolling Hash} ($2^{61}-1$ Mersenne prime).
    \end{itemize}
    \item \textbf{Periodicity \& Border Properties}:
    \begin{itemize}
        \item Minimal period of length-$n$ string is $p = n - \pi[n - 1]$ (if $n \bmod p == 0 \implies S = (S[0..p-1])^{n/p}$).
        \item \textbf{Fine-Wilf (Periodicity Lemma)}: If $S$ has periods $p, q$ and $|S| \ge p + q - \gcd(p, q) \implies S$ has period $\gcd(p, q)$.
        \item \textbf{Border Series}: All borders of a string partition into $\le O(\log n)$ arithmetic progressions $\implies$ Optimize border DP to $O(n \log n)$.
    \end{itemize}
    \item \textbf{Theoretical Bounds}:
    \begin{itemize}
        \item Distinct palindromic substrings in length-$n$ string is always $\le n$.
        \item Suffix Automaton of length-$n$ string has $\le 2n - 1$ states and $\le 3n - 4$ transition edges.
    \end{itemize}
\end{itemize}

\Algorithm{Rolling Hash (61-bit Mersenne)}
{Double/61-bit Mersenne prime Hash ($2^{61}-1$) with randomized base}
{$O(n)$ build, $O(1)$ substring query}
{}{source/pseudocode/rolling_hash.txt}

\Algorithm{KMP}
{Prefix function $\pi$ and linear string matching}
{$O(n + m)$}
{}{source/pseudocode/kmp.txt}

\Algorithm{Z-Algorithm}
{Z-array: $z[i] = \text{LCP}(S, S[i..n-1])$ and pattern matching via $S = P + \# + T$}
{$O(n + m)$}
{}{source/pseudocode/z_algorithm.txt}

\Algorithm{Aho Corasick}
{Multi-pattern dictionary matching on Trie with BFS failure/exit links}
{$O(\sum |P_i| + |T|)$}
{}{source/pseudocode/aho_corasick.txt}

\Algorithm{Suffix Array}
{Suffix Array doubling construction and Kasai LCP algorithm}
{$O(n \log n)$ construction, $O(n)$ LCP}
{}{source/pseudocode/suffix_array.txt}

\begin{itemize}
    \item \textbf{LCP Applications \& Core CP Recipes}:
    \begin{itemize}
        \item \textbf{Arbitrary Suffix LCP}: $\text{LCP}(S[i..], S[j..]) = \min_{k=\min(pos[i], pos[j])+1}^{\max(pos[i], pos[j])} LCP[k]$ via RMQ / Sparse Table in $O(1)$ query.
        \item \textbf{Distinct Substrings}: Total distinct substrings is $\frac{n(n+1)}{2} - \sum_{i=1}^{n-1} LCP[i]$.
        \item \textbf{$k$-th Lexicographical Substring}: Suffix $SA[i]$ contributes $n - SA[i] - LCP[i]$ new strings of lengths $LCP[i]+1 \dots n - SA[i]$.
        \item \textbf{LCS of 2 Strings}: Build SA on $S_1 + \# + S_2$; ans is $\max LCP[i]$ where $SA[i], SA[i-1]$ belong to different strings.
        \item \textbf{LCS of $K$ Strings}: Build SA on $S_1 + \#_1 + \dots + S_K + \#_K$; two pointers / sliding window $[L, R]$ covering all $K$ strings $\implies \min_{k=L+1}^R LCP[k]$.
        \item \textbf{Pattern Occurrences Range}: Suffixes with prefix $P$ form continuous range $[L, R]$ in SA; find via 2 binary searches in $O(|P| \log n)$, count $= R - L + 1$.
        \item \textbf{Longest Substring Occurring $\ge K$ Times}: $\max_i (\min_{j=i+1}^{i+K-1} LCP[j])$ via monotonic queue / sliding RMQ in $O(n)$.
    \end{itemize}
\end{itemize}

\Algorithm{Suffix Automaton}
{Minimal DAWG where every path from root corresponds to a substring of $S$}
{$O(n |\Sigma|)$}
{}{source/pseudocode/suffix_automaton.txt}

\Algorithm{Lyndon Factorization}
{Duval algorithm: factorize string into Lyndon words $s_1 \ge s_2 \ge \dots \ge s_k$ \& Minimal string rotation}
{$O(n)$ time, $O(1)$ space}
{}{source/pseudocode/lyndon.txt}

\Algorithm{Manacher}
{Find the longest palindrome radius centered at every index}
{$O(n)$}
{}{source/pseudocode/manacher.txt}

\Algorithm{Palindrome Tree}
{EERTREE: Represent all distinct palindromes in string and transition links}
{$O(n)$}
{}{source/pseudocode/palindrome_tree.txt}

\Algorithm{De Bruijn (FKM)}
{Stream a cyclic sequence containing every length-$n$ word once: concatenate Lyndon words whose lengths divide $n$, in lexicographic order. For a linear form, append $n-1$ cyclic prefix symbols.}
{$O(k^n)$ time, $O(n)$ auxiliary space; amortized $O(1)$ per emitted symbol}
{}{source/pseudocode/debruijn.txt}


\section{Quy hoạch động (DP)}

\Algorithm{Reroot DP}
{Template reroot trên cây: DFS hai lượt, chuyển trạng thái bằng prefix/suffix}
{$\mathcal{O}(n)$ cho lượt DFS}
{compactcpp}{source/cpp/dp/reroot_dp.cpp}

\Algorithm{Divide-and-Conquer DP và bất đẳng thức tứ giác}
{\textbf{Bất đẳng thức tứ giác (Monge)}: $\forall a \le b \le c \le d \implies C(a, c) + C(b, d) \le C(a, d) + C(b, c)$ ($\text{Cross} \le \text{Nest}$).
\begin{itemize}
  \item \textbf{Kiểm tra cục bộ cần và đủ}: Chỉ cần kiểm tra các ô $2 \times 2$ kề nhau (tổng rút gọn):
  $$ C(i, j+1) - C(i, j) \ge C(i+1, j+1) - C(i+1, j) \quad (\forall i < j) $$
  (Tương đương, độ tăng chi phí $\Delta_j(i) = C(i, j+1) - C(i, j)$ không tăng theo $i$; trong trường hợp liên tục $\frac{\partial^2 C}{\partial x \partial y} \le 0$).
  \item \textbf{Các hàm chi phí thường thỏa QI}:
  \begin{itemize}
    \item $C(l, r) = f(S_r - S_{l-1})$ với $f$ lồi ($f'' \ge 0$), ví dụ $(S_r - S_{l-1})^2$, $(S_r - S_{l-1})^k$ ($k \ge 1$).
    \item $C(l, r) = \sum_{k=l}^r w_k = W_r - W_{l-1}$ ($w_k \ge 0$, dấu bằng trong Knuth Optimal BST).
    \item $C(l, r) = \text{inv}(l, r)$ (số nghịch thế trong đoạn $[l, r]$, hoặc tính chất trên cặp chéo).
    \item $C(l, r) = |x_r - x_l|^p$ ($p \ge 1$ với $x$ đã sắp xếp), hoặc phương sai/tổng bình phương chênh lệch.
  \end{itemize}
  \item \textbf{Đơn điệu}: Tính chất Monge $\implies opt(i_1) \le opt(i_2)$ với $i_1 \le i_2$ (DnC: $O(k n \log n)$) và $opt[i][j-1] \le opt[i][j] \le opt[i+1][j]$ (Knuth: $O(n^2)$).
\end{itemize}
}
{$\mathcal{O}(k n \log n)$ với DnC / $\mathcal{O}(n^2)$ với Knuth}
{compactcpp}{source/cpp/dp/dnc_dp.cpp}

\Algorithm{Tối ưu DP 1D/1D}
{Tối ưu $dp[i] = \min_{j < i} (dp[j] + cost(j, i))$ thỏa bất đẳng thức tứ giác bằng deque các khoảng tối ưu}
{$\mathcal{O}(n \log n)$}
{compactcpp}{source/cpp/dp/one_d_one_d_dp.cpp}

\Algorithm{Alien DP}
{DP trên hàm lồi/lõm (WQS Binary Search):
\begin{itemize}
  \item Chi phí lồi/lõm $\to$ Chọn greedy.
  \item DP lồi/lõm $\to$ Binary search hệ số phạt $\lambda$.
\end{itemize}
}
{$\mathcal{O}(n \log (\text{max\_val}))$}
{compactcpp}{source/cpp/dp/alien_dp.cpp}

\Algorithm{Convex Hull Trick}
{Tối ưu DP với hàm tuyến tính $y = mx + c$}
{$\mathcal{O}(n \log n)$ hoặc $\mathcal{O}(n)$ khi hệ số góc đơn điệu}
{compactcpp}{source/cpp/dp/convex_hull_trick.cpp}

\Algorithm{Slope Trick}
{Duy trì hàm lồi tuyến tính từng đoạn liên tục $f(x)$: cộng $|x-a|$, prefix min, dịch miền bằng hai heap}
{$\mathcal{O}(n \log n)$}
{compactcpp}{source/cpp/dp/slope_trick.cpp}

\subsection{Bitmask DP}
\begin{itemize}
  \item \textbf{Matching}: Tập cặp không có thứ tự $\to$ cố định một đầu, chọn đầu còn lại. $O(n2^n)$ thời gian, $O(2^n)$ bộ nhớ.
  \item \textbf{Nén trạng thái}: $i=\operatorname{popcount}(mask)$ $\to$ $dp[i][mask]\to dp[mask]$.
\end{itemize}

\Algorithm{SOS DP}
{Sum Over Subsets / Fast Zeta Transform: tính tổng trên submask/supermask trong $O(N \cdot 2^N)$ ($N \le 24$)}
{$\mathcal{O}(N \cdot 2^N)$}
{compactcpp}{source/cpp/dp/sos_dp.cpp}

\Algorithm{Tối ưu bằng monotonic queue}
{Deque chỉ số tăng dần, tìm min trong sliding window}
{$O(1)$ khấu hao mỗi push/pop}
{compactcpp}{source/cpp/dp/monotonic_queue.cpp}

\Algorithm{Expected Value Trick}
{Chọn tối ưu $k$ trong $n$ phần tử đã xáo trộn: bước $i$ dùng heuristic dải ngẫu nhiên (có thể bỏ lỡ tối ưu), chỉ xét trạng thái trong $\left[\frac{i \cdot k}{n} - C\sqrt{n}, \frac{i \cdot k}{n} + C\sqrt{n}\right]$.}
{$\mathcal{O}(n \sqrt{n})$}
{compactcpp}{source/cpp/dp/expected_value_dp.cpp}


\section{Toán học}

\Algorithm{Công thức số học}
{Số học modulo, rút gọn số mũ Euler mở rộng, Legendre, Wilson, Frobenius Coin, Floor Trick, CRT, Fibonacci}
{}
{}{}

\begin{itemize}
    \item \textbf{Rút gọn số mũ Euler mở rộng} ($\forall a, m$):
    $$ a^b \equiv a^{(b \bmod \phi(m)) + \phi(m)} \pmod m \quad (b \ge \phi(m)) $$
    \item \textbf{Công thức Legendre ($v_p(n!)$)}: $v_p(n!) = \sum_{k=1}^\infty \lfloor \frac{n}{p^k} \rfloor = \frac{n - S_p(n)}{p - 1}$.
    \item \textbf{LTE}: $v_p(x)$ là số mũ của nguyên tố $p$ trong $x\ne0$. Với $a,b,n$ nguyên dương, $a\ne b$, $n\ge1$: nếu $p$ lẻ, $p\mid a-b$, $p\nmid ab$, thì $v_p(a^n-b^n)=v_p(a-b)+v_p(n)$. Nếu $p$ lẻ, $p\mid a+b$, $p\nmid ab$, $n$ lẻ: $v_p(a^n+b^n)=v_p(a+b)+v_p(n)$. Với $a,b$ lẻ và $p=2$: $n$ lẻ cho $v_2(a^n-b^n)=v_2(a-b)$; $n$ chẵn cho $v_2(a^n-b^n)=v_2(a-b)+v_2(a+b)+v_2(n)-1$. Dùng tính mức chia hết của hiệu/tổng lũy thừa mà không tính số lớn.
    \item \textbf{Định lý Wilson}: $(p - 1)! \equiv -1 \pmod p \iff p$ nguyên tố $\implies (p - 2)! \equiv 1 \pmod p$.
    \item \textbf{Frobenius Coin (Chicken McNugget)}: Với $\gcd(a,b)=1$, số lớn nhất không biểu diễn được dạng $ax+by$ ($x,y\ge0$) là $ab-a-b$; số lượng không biểu diễn được là $\frac{(a-1)(b-1)}{2}$.
    \item \textbf{Floor Trick ($O(\sqrt{n})$ nhóm)}:
\begin{lstlisting}
for (int l = 1, r; l <= n; l = r + 1) {
    int v = n / l; r = n / v;
}
\end{lstlisting}
    \item \textbf{Định lý số dư Trung Hoa (CRT)}: $x \equiv a_i \pmod{m_i}$.
    $$ x \equiv \sum_{i=1}^k a_i M_i (M_i^{-1} \bmod m_i) \pmod M \quad (M = \prod m_i, M_i = M/m_i) $$
    \item \textbf{Fibonacci}: $\gcd(F_m, F_n) = F_{\gcd(m, n)}$, $F_{2n} = F_n(2F_{n+1} - F_n)$.
\end{itemize}

\Algorithm{Like-Euclid Floor Sum}
{Tính $\sum_{i=0}^{n-1} \lfloor \frac{ai+b}{m} \rfloor$ bằng đổi tọa độ}
{$O(\log m)$}
{compactcpp}{source/cpp/math/floor_sum.cpp}

\Algorithm{Rabin-Miller}
{Miller-Rabin tất định: 3 cơ sở khi $n<2^{32}$, 7 cơ sở khi $n<2^{63}$; dùng \texttt{int long long}}
{${O}(k \log n)$}
{compactcpp}{source/cpp/math/rabin_miller.cpp}

\Algorithm{Sàng phân đoạn (Segmented Sieve)}
{Lấy tối đa $k$ nguyên tố trong $[L,R]$, $R\le10^9$; $W=R-L+1$, $B$ số/block; bộ nhớ $O(\sqrt R+B+k)$}
{$O(\sqrt R\log\log R+W\log\log R+\lceil W/B\rceil\pi(\sqrt R))$}
{compactcpp}{source/cpp/math/fast_sieve.cpp}

\Algorithm{Đếm nguyên tố (Prime Counting)}
{Sàng theo các giá trị thương: đếm nguyên tố $\le N$}
{$O(N^{3/4})$ thời gian, $O(\sqrt N)$ bộ nhớ}
{compactcpp}{source/cpp/math/prime_counter.cpp}

\Algorithm{Pollard Rho}
{Tìm ước không tầm thường của hợp số bằng phát hiện chu trình Brent}
{$O(\sqrt[4]{n} \log n)$}
{compactcpp}{source/cpp/math/pollard_rho.cpp}

\Algorithm{Primitive Root}
{Tìm căn nguyên thủy $g$ nhỏ nhất modulo nguyên tố $p$}
{$\mathcal{O}(\text{ans} \cdot \log p)$}
{compactcpp}{source/cpp/math/primitive_root.cpp}

\Algorithm{Baby-step Giant-Step}
{Tìm $x$ nhỏ nhất thỏa $a^x \equiv b \pmod m$ (Discrete Log)}
{$O(\sqrt{m} \log m)$}
{compactcpp}{source/cpp/math/discrete_log_bsgs.cpp}

\Algorithm{Căn modulo (Tonelli-Shanks)}
{Tonelli-Shanks: giải $x^2 \equiv a \pmod p$}
{$O(\log^2 p)$}
{compactcpp}{source/cpp/math/sqrt_mod.cpp}

\noindent\textbf{Bézout (tổng quát)}: Với $a_1,\ldots,a_n\in\mathbb Z$ không đồng thời bằng 0, $g=\gcd(|a_1|,\ldots,|a_n|)$: tồn tại $c_i\in\mathbb Z$ sao cho $\sum_i c_i a_i=g$. Mọi tổ hợp nguyên tạo đúng các bội của $g$: $\sum_i c_i a_i=x$ có nghiệm iff $g\mid x$. \textbf{Cue}: cộng/trừ các bước tùy ý $\to$ gcd các bước chia hết độ dịch chuyển; nghịch đảo modulo $\to ax+my=1$ (cần $\gcd(a,m)=1$). Hệ số được âm; yêu cầu $c_i\ge0$ thì gcd chia hết chưa đủ. Extended GCD tìm hệ số cho hai số, ghép dần cho nhiều số.

\noindent\textbf{GCD: cộng/trừ và đổi mốc}
\begin{itemize}
    \setlength{\itemsep}{1pt}
    \setlength{\parsep}{0pt}
    \item $\gcd(u,v)=\gcd(u,v+ku)=\gcd(v,u\bmod v)$ ($k\in\mathbb Z$, $v\ne0$ cho modulo); $\gcd(cu,cv)=|c|\gcd(u,v)$. Cộng/trừ bội giữ các ước chung.
    \item Đặt $d_i=a_{i+1}-a_i$, $g=\gcd|d_i|$: $\gcd(a_1,\ldots,a_n)=\gcd(a_1,g)$; tịnh tiến: $\gcd(a_1+b,\ldots,a_n+b)=\gcd(a_1+b,g)$ (Row GCD). Phải giữ mốc; $\gcd(z,0)=|z|$.
    \item $\gcd(pu+qv,ru+sv)=\gcd(u,v)$ nếu $ps-qr=\pm1$ (phép ngược có hệ số nguyên). Tổ hợp tùy ý chưa chắc giữ gcd, vd. $(u+v,u-v)$.
\end{itemize}

\Algorithm{Extended GCD}
{Tìm $x,y$ thỏa phương trình $ax + by = \gcd(a, b)$}
{$\mathcal{O}(\log \min(a,b))$}
{compactcpp}{source/cpp/math/extended_gcd.cpp}

\Algorithm{Diophantine}
{Nghiệm tổng quát của phương trình Diophantine tuyến tính $ax + by = c$}
{$O(\log \min(a, b))$}
{compactcpp}{source/cpp/math/diophantine.cpp}

\Algorithm{Stern-Brocot Tree}
{Phân số tối giản, mã hóa/giải mã cây, LCA}
{$\mathcal{O}(\log(p + q))$}
{compactcpp}{source/cpp/math/stern_brocot.cpp}

\Algorithm{Hàm nhân tính và M\"obius Inversion}
{Tích chập Dirichlet, nghịch đảo M\"obius, tổng GCD/LCM}
{}
{}{}

\begin{itemize}
    \item \textbf{Euler Phi \& M\"obius}: $\phi(n) = n \prod_{p \mid n} (1 - 1/p)$, $\sum_{d \mid n} \phi(d) = n$, $\sum_{d \mid n} \mu(d) = [n = 1]$.
    \item \textbf{Dirichlet \& Inversion}: $(f * g)(n) = \sum_{d \mid n} f(d) g(n / d)$.
    $$ f(n) = \sum_{d \mid n} g(d) \iff g(n) = \sum_{d \mid n} \mu(d) f(n / d) $$
    \item \textbf{Đếm cặp nguyên tố cùng nhau}:
    $$ \sum_{i=1}^n \sum_{j=1}^m [\gcd(i, j) = 1] = \sum_{d=1}^{\min(n, m)} \mu(d) \left\lfloor \frac{n}{d} \right\rfloor \left\lfloor \frac{m}{d} \right\rfloor $$
    \item \textbf{Tổng GCD trên mọi cặp}:
    $$ \sum_{i=1}^n \sum_{j=1}^m \gcd(i, j) = \sum_{g=1}^{\min(n, m)} \phi(g) \left\lfloor \frac{n}{g} \right\rfloor \left\lfloor \frac{m}{g} \right\rfloor $$
    \item \textbf{Tổng LCM trên mọi cặp} ($S(x) = \frac{x(x+1)}{2}$):
    $$ \sum_{i \le n, j \le m} \text{lcm}(i, j) = \sum_{g} g \sum_{d} d^2 \mu(d) S\left(\left\lfloor \frac{n}{gd} \right\rfloor\right) S\left(\left\lfloor \frac{m}{gd} \right\rfloor\right) $$
\end{itemize}

\Algorithm{Công thức tổ hợp}
{Đẳng thức nhị thức, Hockey-stick, Vandermonde, Inversion, Stars and Bars, Catalan, Stirling, Cayley, LGV}
{}
{}{}

\begin{itemize}
    \item \textbf{Đẳng thức nhị thức}: $\binom{n}{k} = \binom{n}{n-k} = \binom{n-1}{k-1} + \binom{n-1}{k}$; $k \binom{n}{k} = n \binom{n-1}{k-1}$; $\frac{1}{k+1}\binom{n}{k} = \frac{1}{n+1}\binom{n+1}{k+1}$; $\binom{n}{k} \binom{k}{m} = \binom{n}{m} \binom{n-m}{k-m}$; Đối số trên âm: $\binom{-n}{k} = (-1)^k \binom{n+k-1}{k}$.
    \item \textbf{Tổng nhị thức}: $\sum_{k=0}^n \binom{n}{k} = 2^n$; $\sum_{k=0}^n (-1)^k \binom{n}{k} = [n = 0]$; $\sum_{k \text{ chẵn}} \binom{n}{k} = \sum_{k \text{ lẻ}} \binom{n}{k} = 2^{n-1}$; $\sum_{k=0}^n k \binom{n}{k} = n 2^{n-1}$; $\sum_{k=0}^n k^2 \binom{n}{k} = n(n+1) 2^{n-2}$; $\sum_{k=0}^n \frac{1}{k+1}\binom{n}{k} = \frac{2^{n+1}-1}{n+1}$.
    \item \textbf{Hockey-Stick \& Vandermonde}: $\sum_{i=r}^n \binom{i}{r} = \binom{n+1}{r+1}$; $\sum_{i=0}^k \binom{n+i}{i} = \binom{n+k+1}{k}$; $\sum_{k=0}^r \binom{m}{k} \binom{n}{r-k} = \binom{m+n}{r} \implies \sum_{k=0}^n \binom{n}{k}^2 = \binom{2n}{n}$.
    \item \textbf{Binomial Inversion và sai phân}: $f(n) = \sum_{k=0}^n \binom{n}{k} g(k) \iff g(n) = \sum_{k=0}^n (-1)^{n-k} \binom{n}{k} f(k)$; Dạng tổng phía trên: $f(n) = \sum_{k=n}^N \binom{k}{n} g(k) \iff g(n) = \sum_{k=n}^N (-1)^{k-n} \binom{k}{n} f(k)$; $\sum_{k=0}^n (-1)^{n-k} \binom{n}{k} k^m = n! \left\{ \begin{matrix} m \\ n \end{matrix} \right\}$ ($0$ nếu $m<n$, $n!$ nếu $m=n$).
    \item \textbf{Stars and Bars}: $x_i \ge 0 \implies \binom{n + k - 1}{k - 1}$; $x_i \ge 1 \implies \binom{n - 1}{k - 1}$; $x_i \ge c_i \implies n' = n - \sum c_i$.
    \item \textbf{Derangements}: $D_0 = 1, D_1 = 0, D_n = (n - 1)(D_{n-1} + D_{n-2}) = n D_{n-1} + (-1)^n \approx [n! / e]$.
    \item \textbf{Số Catalan $C_n = \frac{1}{n+1}\binom{2n}{n}$}: Dãy ngoặc đúng, cây nhị phân đầy đủ có $n+1$ lá, chia tam giác đa giác lồi $(n+2)$ đỉnh, đường Dyck không vượt $y=x$, hoán vị sinh bằng một stack.
    \item \textbf{Nguyên lý phản xạ (Reflection Principle)}: Đường đi $(0,0)\to(n,m)$ không chạm $y = x + k$ ($k > 0$): $\binom{n + m}{m} - \binom{n + m}{m - k}$.
    \item \textbf{Định lý Ballot}: A có $a$ phiếu, B có $b$ ($a>b$). Xác suất A luôn dẫn trước nghiêm ngặt: $\frac{a - b}{a + b}$.
    \item \textbf{Số Stirling loại 1 và 2}: $\left[ \begin{matrix} n \\ k \end{matrix} \right] = \left[ \begin{matrix} n-1 \\ k-1 \end{matrix} \right] + (n-1) \left[ \begin{matrix} n-1 \\ k \end{matrix} \right]$, $\left\{ \begin{matrix} n \\ k \end{matrix} \right\} = \left\{ \begin{matrix} n-1 \\ k-1 \end{matrix} \right\} + k \left\{ \begin{matrix} n-1 \\ k \end{matrix} \right\}$.
    $x^n = \sum_{k=0}^n \left\{ \begin{matrix} n \\ k \end{matrix} \right\} x^{\underline{k}} \implies \sum_{i=1}^n i^k = \sum_{j=0}^k \left\{ \begin{matrix} k \\ j \end{matrix} \right\} j! \binom{n + 1}{j + 1}$.
    \item \textbf{Chu kỳ hoán vị}: Với hoán vị $\pi$ của $n$ phần tử có $c$ chu kỳ (tính cả điểm cố định), số phép đổi chỗ hai vị trí bất kỳ ít nhất để về identity là $n-c$. Mỗi chu kỳ dài $\ell$ cần $\ell-1$ lần. Biến hoán vị $A$ thành $B$: xét $\pi(i)=\operatorname{pos}_B(A_i)$. Không áp dụng cho chỉ đổi hai vị trí kề nhau.
    \item \textbf{Eulerian numbers}: $E(n,k)$ đếm hoán vị $n$ phần tử có đúng $k$ vị trí giảm $a_i>a_{i+1}$. $E(0,0)=1$; với $n\ge1$, $E(n,k)=(n-k)E(n-1,k-1)+(k+1)E(n-1,k)$, chỉ tính $0\le k<n$, các chỉ số khác bằng $0$. Với $n\ge1$, cũng bằng số hoán vị có đúng $k$ vị trí $a_i>i$ (chỉ số từ $1$).
    \item \textbf{PIE}: $|\bigcup_{i=1}^n A_i| = \sum_{\emptyset \ne J} (-1)^{|J|-1} |\bigcap_{j \in J} A_j|$; Thuộc đúng $r$ tập: $\sum_{m=r}^n (-1)^{m-r} \binom{m}{r} \sum_{|X|=m} |\bigcap_{i \in X} A_i|$.
    \item \textbf{Định lý Lucas}: $\binom{n}{k} \equiv \prod_{i} \binom{n_i}{k_i} \pmod p$ (chữ số cơ số $p$).
    \item \textbf{Đếm cây và Cayley}: $n^{n-2}$ cây có nhãn, rừng gồm $k$ cây có gốc: $k n^{n-k-1}$.
    \item \textbf{LGV (Lindstr\"om--Gessel--Viennot)}: Trên DAG, đặt $M_{ij}$ là số đường đi $A_i\to B_j$. Khi đó $\det M=\sum_{\sigma\in S_k}\operatorname{sgn}(\sigma)N_\sigma$, với $N_\sigma$ đếm các bộ đường đi đôi một không chung đỉnh, nối $A_i\to B_{\sigma(i)}$. Nếu mọi phép ghép khác đồng nhất đều buộc giao nhau, $\det M$ chính là số bộ nối $A_i\to B_i$. Nhận biết: nhiều đường đi đơn điệu trên lưới, không chung đỉnh; mỗi $M_{ij}=\binom{\Delta x+\Delta y}{\Delta x}$ nếu $\Delta x,\Delta y\ge0$, ngược lại $0$. Các bộ giao nhau triệt tiêu bằng cách đổi đuôi sau điểm giao, làm đổi dấu hoán vị.
    \item \textbf{Matrix--Tree (Kirchhoff)}: Đồ thị vô hướng có Laplacian $L_{ii}=\deg(i)$, $L_{ij}=-m_{ij}$ ($i\ne j$), với $m_{ij}$ là số cạnh nối $i,j$. Xóa hàng và cột cùng chỉ số $r$ được $L^{(r)}$; số cây khung $\tau(G)=\det L^{(r)}$, với mọi $r$. Đồ thị không liên thông cho $0$; một đỉnh cho $1$. Giữ cạnh song song, bỏ self-loop. Tính modulo nguyên tố bằng \texttt{determinant\_mod}.
    \item \textbf{Đếm MST}: Kruskal theo nhóm cạnh cùng trọng số; co DSU từ các cạnh nhẹ hơn, bỏ cạnh nội bộ. Nhân $\tau(C)$ cho từng thành phần liên thông $C$ của đồ thị co (giữ cạnh song song), rồi mới union nhóm hiện tại. Đồ thị gốc không liên thông: không có MST.
\end{itemize}

\Algorithm{Hàm sinh (Generating Functions) và chuỗi Taylor}
{OGF, EGF, khai triển Taylor, tích phân số Simpson}
{}
{}{}

\begin{itemize}
    \item OGF: $A(x) = \sum a_n x^n$; EGF: $A(x) = \sum a_n \frac{x^n}{n!}$.
    \item $\frac{1}{1-x} = \sum_{n \ge 0} x^n$, $\frac{1}{(1-x)^k} = \sum_{n \ge 0} \binom{n+k-1}{k-1} x^n$.
    \item $e^x = \sum_{n \ge 0} \frac{x^n}{n!}$, $\ln(1 - x) = -\sum_{n \ge 1} \frac{x^n}{n}$, $(1+x)^\alpha = \sum_{n \ge 0} \binom{\alpha}{n} x^n$.
    \item \textbf{Công thức Simpson}: $\int_a^{a+2h} f(x)dx \approx \frac{h}{3}(f(a) + 4f(a+h) + f(a + 2h))$.
\end{itemize}

\Algorithm{Truy hồi tuyến tính (Berlekamp-Massey)}
{Tìm truy hồi tuyến tính tối tiểu từ $2N$ số hạng đầu, tính số hạng thứ $K$ qua $x^K\bmod P(x)$}
{${O}(N^2)$ tìm truy hồi, $\mathcal{O}(N^2 \log K)$ tính số hạng thứ $k$}
{compactcpp}{source/cpp/math/berlekamp_massey.cpp}

\Algorithm{FWHT}
{Fast Walsh-Hadamard Transform (XOR, AND, OR Convolution)}
{$\mathcal{O}(n \log n)$}
{compactcpp}{source/cpp/math/fwht.cpp}

\Algorithm{Lagrange}
{Nội suy Lagrange tại $X$ từ $n+1$ điểm cách đều có giá trị $y_0,\dots,y_n$}
{$O(n)$ tính giá trị}
{compactcpp}{source/cpp/math/lagrange.cpp}
\textbf{Công thức:} Với $d+1$ điểm phân biệt $(x_i,y_i)$ và $\deg f\le d$,
\[
f(X)=\sum_{i=0}^{d}y_i\prod_{\substack{0\le j\le d\\j\ne i}}\frac{X-x_j}{x_i-x_j}.
\]
Code dùng $x_i=i$, nên $\prod_{j\ne i}(i-j)=(-1)^{d-i}i!(d-i)!$:
\[
f(X)=\sum_{i=0}^{d}y_i\frac{(-1)^{d-i}}{i!(d-i)!}
\prod_{\substack{0\le j\le d\\j\ne i}}(X-j).
\]
Trong modulo $p$, phép chia là nhân nghịch đảo; tích tiền tố/hậu tố tính các tích bỏ $i$ trong $O(d)$.\\
\textbf{Khi dùng:} Cần $f(N)$ với $N$ rất lớn, biết $f$ là đa thức bậc tối đa $d$ và tính được $f(0),\dots,f(d)$ nhanh. Không suy ra tính đa thức chỉ từ vài giá trị đầu.\\
\textbf{Ví dụ:} $S_k(N)=\sum_{i=1}^{N}i^k$ có bậc $k+1$: lập $y[0..k+1]$ bằng tổng tiền tố, gọi \texttt{lagrange(y,N,p)}. Tương tự, đếm cách chọn $k$ phần tử từ $N$ phần tử: $\binom{N}{k}$ là đa thức bậc $k$.\\
\textbf{Điều kiện code:} $y[i]=f(i)$, modulo nguyên tố $p$, số điểm $d+1<p$; với tổng lũy thừa cần $k+2<p$. Mỗi lần tính tốn $O(d+\log p)$, gồm nghịch đảo modulo.

\Algorithm{NTT (Number Theoretic Transform)}
{NTT lặp với bit-reversal, MOD = $998244353$ ($G=3$)}
{$\mathcal{O}(n \log n)$}
{compactcpp}{source/cpp/math/ntt.cpp}

\Algorithm{FFT (Fast Fourier Transform)}
{FFT Cooley-Tukey trên số phức cho tích chập số thực/tổng quát}
{$\mathcal{O}(n \log n)$}
{compactcpp}{source/cpp/math/fft.cpp}

\Algorithm{Công thức Formal Power Series (FPS)}
{Công thức thao tác đa thức bằng phương pháp Newton và convolution}
{}
{}{}

\begin{itemize}
    \item \textbf{Nghịch đảo ($Q=P^{-1}$)}: $Q_k \equiv Q_{k-1}(2 - P Q_{k-1}) \pmod{x^{2^k}}$, khởi đầu $Q_0 = P(0)^{-1}$.
    \item \textbf{Đạo hàm và tích phân}: $P'(x) = \sum i a_i x^{i-1}$, $\int P(x) = \sum \frac{a_i}{i+1} x^{i+1}$.
    \item \textbf{Logarithm}: $\ln P(x) = \int \frac{P'(x)}{P(x)} dx \pmod{x^n}$ (cần $P(0) = 1$).
    \item \textbf{Hàm mũ}: $Q_k \equiv Q_{k-1}(1 - \ln Q_{k-1} + P) \pmod{x^{2^k}}$ (cần $P(0) = 0$).
    \item \textbf{Căn bậc hai}: $Q_k \equiv \frac{1}{2}(Q_{k-1} + P Q_{k-1}^{-1}) \pmod{x^{2^k}}$ ($P(0)$ là thặng dư bậc hai).
    \item \textbf{Lũy thừa}: $P(x)^k = \exp(k \ln P(x))$. Nếu $P(x) = c x^d (1 + Q(x))$, $P(x)^k = c^k x^{k d} \exp(k \ln(1 + Q(x)))$.
    \item \textbf{Chia lấy thương và dư ($A = BQ + R$)}: Với $\text{rev}_k(P) = x^k P(1/x)$:
    $$\text{rev}_{n-m}(Q) \equiv \text{rev}_n(A) \cdot (\text{rev}_m(B))^{-1} \pmod{x^{n-m+1}}, \quad R = A - BQ$$
    \item \textbf{Taylor Shift $P(x + c)$}: $P(x + c) = \sum_{k=0}^n \frac{x^k}{k!} \sum_{i=k}^n (i! a_i) \frac{c^{i-k}}{(i-k)!}$. Convolution dãy $A_i=i!a_i$ đảo ngược với $B_j = \frac{c^j}{j!}$.
    \item \textbf{Convolution modulo tùy ý (MTT)}: Tách $A=A_1B+A_0$, $B= \lceil\sqrt{M}\rceil \approx 32000$. Tính $(A_1B_1),(A_1B_0+A_0B_1),(A_0B_0)$ bằng bốn FFT.
\end{itemize}

\Algorithm{Đại số tuyến tính (Gauss)}
{Khử Gauss modulo nguyên tố $p$: định thức, hạng, nghiệm riêng và cơ sở nghiệm thuần nhất}
{$\mathcal{O}(n^3)$}
{compactcpp}{source/cpp/math/linalg_gauss.cpp}
Nghiệm tổng quát: $x=x_0+\sum_i t_i v_i\pmod p$, với $t_i$ tùy ý; 	exttt{consistent=false} nếu vô nghiệm.

\Algorithm{XOR Linear Basis}
{XOR basis: XOR max/min/thứ $k$ của tập con, sinh $2^{\text{size}}$ giá trị phân biệt}
{$O(B)$ thêm/truy vấn ($B \le 60$)}
{compactcpp}{source/cpp/math/xor_basis.cpp}

\Algorithm{Burnside Lemma \& P\'olya Theorem}
{Đếm lớp tương đương dưới tác động nhóm: necklace, bracelet}
{}
{}{}

\begin{itemize}
    \item $|Classes| = \frac{1}{|G|} \sum_{g \in G} |Fix(g)| = \frac{1}{|G|} \sum_{g \in G} k^{c(g)}$.
    \item \textbf{Necklace (chỉ $n$ phép xoay)}: $N(n, k) = \frac{1}{n} \sum_{d \mid n} \phi(d) k^{n / d}$.
    \item \textbf{Bracelet (xoay và đối xứng)}:
    \begin{itemize}
        \item $n$ lẻ: $B(n, k) = \frac{1}{2} N(n, k) + \frac{1}{2} k^{(n+1)/2}$.
        \item $n$ chẵn: $B(n, k) = \frac{1}{2} N(n, k) + \frac{1}{4} (k^{n/2} + k^{n/2 + 1})$.
    \end{itemize}
\end{itemize}

\subsection{Game: Invariant, Symmetric, Nim, DP}
% References: MIT ES.268 Theory of Impartial Games; CMU Combinatorics F11/ex4a.pdf.
\begin{itemize}
    \setlength{\itemsep}{0pt}
    \setlength{\parsep}{0pt}
    \item \textbf{1. Invariant (Bất biến)}: Tìm đại lượng bất biến (parity chẵn/lẻ, modulo $k+1$, số nghịch thế) mà mọi nước đi hợp lệ đều bảo toàn hoặc đổi luân phiên. Người thắng luôn đưa trạng thái về tập bất biến mà đối thủ không thể duy trì. \textit{Ví dụ}: Bốc $1\dots k$ sỏi từ $n$ viên (lấy cuối thắng): Bất biến $n\equiv 0\pmod{k+1}$ là thế thua (P-pos). Nếu $n\not\equiv 0$, bốc $n\bmod(k+1)$ về bội $k+1$; đối thủ bốc $x\in[1,k]$, ta bốc $k+1-x$ để giữ $0\pmod{k+1}$.
    \item \textbf{2. Symmetric (Đối xứng / Bắt cặp)}: Ghép các trạng thái hoặc nước đi thành cặp đối xứng rời nhau (qua tâm, trục, hoặc perfect matching). Đối thủ đi vào một nửa thì ta đi vào nửa còn lại để khôi phục đối xứng. Người đầu thắng nếu chiếm tâm/nước lẻ đầu rồi phản chiếu; người sau thắng nếu bàn cờ ban đầu đã hoàn toàn đối xứng. \textit{Ví dụ}: Đặt xu lên bàn tròn: P1 đặt đúng tâm; P2 đặt tại $(x,y)$ thì P1 đặt tại $(-x,-y)\implies$ P1 thắng. Trên đồ thị có perfect matching: P2 luôn đi theo cạnh ghép của đỉnh P1 vừa chọn $\implies$ P2 thắng.
    \item \textbf{3. Nim (XOR \& Sprague-Grundy)}: Impartial game dưới normal play (hết nước đi là thua). Mọi game trên DAG tương đương 1 đống Nim kích thước $G(u)=\text{mex}\{G(v)\mid u\to v\}$. Độc lập: $G_{\text{tổng}}=\bigoplus_i G_i$. Thua (P-pos) $\iff G=0$; Thắng (N-pos) $\iff G\ne 0$ (chọn đống có bit cao nhất của $G$, giảm $a_i\to a_i\oplus G$).
    \emph{Mis\`ere Nim} (lấy cuối thua): mọi đống $\le1\implies$ thua $\iff$ số đống 1 lẻ; có đống $>1\implies$ chơi như Nim thường, còn đúng 1 đống $>1$ thì giảm nó về $0$ hoặc $1$ để lại số đống 1 lẻ.
    \emph{Staircase Nim}: chuyển sỏi từ bậc $i\to i-1$ (bậc 0 loại bỏ); chỉ XOR các bậc \emph{lẻ} ($G=\bigoplus_{i\text{ lẻ}} a_i$), đối thủ dời chẵn $\to$ lẻ thì ta dời tiếp số đó lẻ $\to$ chẵn kề dưới.
    \emph{Wythoff Nim}: bốc 1 đống tùy ý hoặc 2 đống cùng lượng; thua $\iff a=\lfloor k\varphi\rfloor, b=a+k$ với $k=b-a, \varphi=(1+\sqrt5)/2$.
    \emph{Green Hackenbush (cây)}: chặt cạnh, nhánh rời gốc rơi: $g(u)=\bigoplus_{v\in\text{child}(u)}(g(v)+1)$.
    \item \textbf{4. DP (Quy hoạch động \& Retrograde)}: Biểu diễn trạng thái trên đồ thị có hướng. Trạng thái lá (hết nước đi) là \textbf{LOSE}. Chuyển trạng thái: $\exists v\in\text{LOSE}\implies u\in\text{WIN}$; $\forall v\in\text{WIN}\implies u\in\text{LOSE}$.
    \emph{Retrograde (chu trình/hòa)}: BFS ngược từ lá (LOSE), theo dõi bậc ra $\text{out\_deg}[u]$: kề LOSE $\implies$ WIN; kề WIN $\implies$ giảm $\text{out\_deg}[u]$, về $0\implies$ LOSE; còn lại $\implies$ \textbf{DRAW} ($O(V+E)$). \textit{Ví dụ}: Game chia kẹo/trừ số: $dp[i]=\text{LOSE}\iff\forall\text{step}, dp[i-\text{step}]=\text{WIN}$. Game di chuyển quân cờ lặp trạng thái: dùng Retrograde queue.
\end{itemize}

\Algorithm{Xác suất và kỳ vọng}
{Tính tuyến tính của kỳ vọng, phân tích biến chỉ báo, Coupon Collector, bất đẳng thức}
{}
{}{}

\begin{itemize}
    \setlength{\itemsep}{0pt}
    \setlength{\parsep}{0pt}
    \item \textbf{Tuyến tính kỳ vọng}: $E[\sum c_i X_i] = \sum c_i E[X_i]$ (luôn đúng $\forall X_i$, \emph{kể cả khi các biến ngẫu nhiên phụ thuộc nhau}).
    \item \textbf{Biến chỉ báo (Indicator Variable Analysis)}: Đặt $I_A = [A\text{ xảy ra}] \in \{0, 1\}$.
    Khi đó $E[I_A] = 1\cdot P(A) + 0\cdot P(\overline{A}) = P(A)$.
    \begin{itemize}
        \setlength{\itemsep}{0pt}
        \setlength{\parsep}{0pt}
        \item \emph{Kỳ vọng đếm}: Muốn tính kỳ vọng số đối tượng thỏa mãn tính chất, viết $X = \sum_{i=1}^n I_{A_i} \implies E[X] = \sum_{i=1}^n P(A_i)$. Ta chuyển việc tính kỳ vọng toàn cục phức tạp thành tính xác suất xảy ra của từng biến cố thành phần riêng lẻ (nghiệm đúng cả khi các $A_i$ phụ thuộc nhau).
        \item \emph{Ví dụ}: (1) Hoán vị ngẫu nhiên: số nghịch thế $E = \binom{n}{2} \cdot \frac{1}{2} = \frac{n(n-1)}{4}$; số điểm bất động $\pi(i)=i$ là $E = n \cdot \frac{1}{n} = 1$. (2) Đồ thị $G(n, p)$: số tam giác $K_3$ có $E = \binom{n}{3}p^3$; số đỉnh cô lập có $E = n(1-p)^{n-1}$.
        \item \emph{Phương sai \& Tích}: $I_A^2 = I_A \implies E[I_A^2] = P(A)$; $I_A I_B = I_{A \cap B} \implies E[I_A I_B] = P(A \cap B)$. Khi đó $\text{Cov}(I_A, I_B) = P(A \cap B) - P(A)P(B)$; $\text{Var}(X) = \sum \text{Var}(I_i) + 2\sum_{i<j} \text{Cov}(I_i, I_j)$.
    \end{itemize}
    \item \textbf{Phân phối hình học}: $E[X] = 1/p$, $\text{Var}(X) = (1-p)/p^2$.
    \item \textbf{Coupon Collector}: $E[T] = n \sum_{i=1}^n \frac{1}{i} = n H_n \approx n(\ln n + \gamma)$ ($\gamma \approx 0.5772$).
    \item \textbf{Bất đẳng thức}: \textbf{Markov}: $P(X \ge a) \le \frac{E[X]}{a}$ ($X \ge 0$); \textbf{Chebyshev}: $P(|X - \mu| \ge k\sigma) \le \frac{1}{k^2}$.
\end{itemize}

\Algorithm{Bignum}
{Các block cơ số $10^9$; tính trung gian bằng 64-bit. In block cao nhất bình thường, các block còn lại đủ 9 chữ số.}
{$O(n)$ in ($n$ block)}
{compactcpp}{source/cpp/math/big_integer.cpp}


\section{Geometry}

\Algorithm{Primitive Geometry}
{Core 2D primitives: Points, Vectors, Orientation, Line Intersections, Projections}
{$\mathcal{O}(1)$ per operation}
{}{source/pseudocode/primitive_geometry.txt}

\Algorithm{Convex Hull}
{Monotone Chain / Graham Scan: Smallest convex polygon enclosing point set}
{$\mathcal{O}(n \log n)$}
{}{source/pseudocode/convex_hull.txt}

\Algorithm{Rotating Calipers}
{Rotating Calipers: Polygon diameter (farthest pair), Minimum bounding rectangle}
{$\mathcal{O}(n)$}
{}{source/pseudocode/rotating_calipers.txt}

\Algorithm{Minkowski Sum}
{Minkowski Sum of two convex polygons. Min polygon distance: $\text{dist}(P, Q) = \text{dist}(O, P - Q)$}
{$\mathcal{O}(|P| + |Q|)$}
{}{source/pseudocode/minkowski_sum.txt}

\Algorithm{Point in Convex Polygon}
{Test if point lies inside convex polygon via polar angle binary search from $P_0$}
{$\mathcal{O}(\log n)$}
{}{source/pseudocode/point_in_convex.txt}

\Algorithm{Line Hull Intersection}
{Binary search extreme vertices along normal vector to find intersection points}
{$O(\log n)$}
{}{source/pseudocode/line_hull_intersection.txt}

\Algorithm{Polygon Raycast}
{Horizontal raycasting point-in-polygon test for arbitrary simple polygons}
{$O(n)$}
{}{source/pseudocode/polygon_raycast.txt}

\Algorithm{Half-plane Intersection}
{Intersection of half-planes forming convex polygon via Deque}
{$O(n \log n)$}
{}{source/pseudocode/halfplane_intersection.txt}

\Algorithm{Smallest Enclosing Circle}
{Welzl's randomized algorithm for minimum bounding circle}
{$\mathcal{O}(n)$}
{}{source/pseudocode/smallest_enclosing_circle.txt}

\Algorithm{Circle Tangents}
{Determine up to 4 common tangent lines between 2 circles}
{$\mathcal{O}(1)$}
{}{source/pseudocode/circle_tangents.txt}

\Algorithm{Geometry Formulas}
{Planar theorems, circles, triangle centers, circle inversion \& 3D primitives}
{}
{}{}

\begin{itemize}
    \item \textbf{Linear Transformation $\leftrightarrow$ Matrix}: Every linear map $T:\mathbb{R}^n\to\mathbb{R}^m$ has $A=[T(e_1)\;\cdots\;T(e_n)]$ such that $T(x)=Ax$. Common 2D forms:
    \begin{itemize}
        \item Rotation by $\theta$: $\begin{pmatrix}\cos\theta&-\sin\theta\\\sin\theta&\cos\theta\end{pmatrix}$; scaling: $\begin{pmatrix}s_x&0\\0&s_y\end{pmatrix}$.
        \item Reflect across $Ox$, $Oy$, $y=x$: $\begin{pmatrix}1&0\\0&-1\end{pmatrix}$, $\begin{pmatrix}-1&0\\0&1\end{pmatrix}$, $\begin{pmatrix}0&1\\1&0\end{pmatrix}$.
        \item Shear: $\begin{pmatrix}1&k\\0&1\end{pmatrix}$ or $\begin{pmatrix}1&0\\k&1\end{pmatrix}$; projection onto unit vector $u$: $P=uu^T$.
        \item Translation is affine, not linear: $x'=Ax+b$. In homogeneous coordinates, $\begin{pmatrix}x'\\y'\\1\end{pmatrix}=\begin{pmatrix}a&b&t_x\\c&d&t_y\\0&0&1\end{pmatrix}\begin{pmatrix}x\\y\\1\end{pmatrix}$.
        \item Composition $T_2\circ T_1$ corresponds to $A_2A_1$; repeating one transform $k$ times corresponds to $A^k$.
    \end{itemize}
    \item \textbf{Pick's Theorem}: $S = I + \frac{B}{2} - 1 \implies I = S - \frac{B}{2} + 1$, where $B = \sum \gcd(|x_{i+1} - x_i|, |y_{i+1} - y_i|)$.
    \item \textbf{Euler's Planar Formula}: $V - E + F = 1 + C$ ($V - E + F = 2$ for connected planar graphs).
    \item \textbf{Ptolemy's Theorem}: Cyclic quad $ABCD \iff AC \cdot BD = AB \cdot CD + BC \cdot AD$.
    \item \textbf{Stewart's Theorem}: $b^2m + c^2n = a(d^2 + mn)$ (cevian $d$ divides side $a$ into $m, n$).
    \item \textbf{Bretschneider's Formula} (General quadrilateral area):
    $$ S = \sqrt{(p-a)(p-b)(p-c)(p-d) - abcd \cos^2\left(\tfrac{A+C}{2}\right)} $$
    \textit{(Cyclic quad $\implies \cos(\frac{A+C}{2}) = 0$, reduces to Brahmagupta's formula)}.
    \item \textbf{Descartes' Theorem} (4 mutually tangent circles):
    $$ \left(\sum_{i=1}^4 k_i\right)^2 = 2 \sum_{i=1}^4 k_i^2 \quad \left(k_i = \pm \frac{1}{R_i}\right) $$
    \item \textbf{Circle Inversion}: Center $O$, radius $R$. $P \mapsto P'$ satisfies $OP \cdot OP' = R^2$. Maps lines not through $O \longleftrightarrow$ circles through $O$; preserves angles and tangencies.
    \item \textbf{Circle-Circle Intersection}: Let $d = |O_1 - O_2|$. Intersects when $|R_1 - R_2| \le d \le R_1 + R_2$, angle $\alpha = \text{acos}\left(\frac{R_1^2 + d^2 - R_2^2}{2 R_1 d}\right)$.
    \item \textbf{Apollonius Circle}: Locus $\frac{PA}{PB} = k$ ($k \ne 1$) is a circle with diameter connecting the internal and external division points of $AB$ with ratio $k$.
    \item \textbf{Circumcenter $O$ \& Orthocenter $H$}: Let $d_B = x_B^2+y_B^2-x_A^2-y_A^2$, $d_C = x_C^2+y_C^2-x_A^2-y_A^2$:
    $$x_O = \frac{d_B (y_C - y_A) - d_C (y_B - y_A)}{D}$$
    $$y_O = \frac{(x_B - x_A) d_C - (x_C - x_A) d_B}{D}$$
    $$D = 2 \left[ (x_B - x_A)(y_C - y_A) - (x_C - x_A)(y_B - y_A) \right]$$
    $$H = A + B + C - 2O$$

    \item \textbf{3D Primitives}:
    \begin{itemize}
        \item 3D Cross Product: $\vec{A} \times \vec{B} = (A_y B_z - A_z B_y, \; A_z B_x - A_x B_z, \; A_x B_y - A_y B_x)$.
        \item Tetrahedron Volume: $V = \frac{1}{6} \left| ((B - A) \times (C - A)) \cdot (D - A) \right|$.
        \item Plane through $P_0$ with normal $\vec{n}$: $\vec{n} \cdot (P - P_0) = 0$; Point-plane distance: $\frac{|\vec{n} \cdot (P - P_0)|}{|\vec{n}|}$.
    \end{itemize}
\end{itemize}

\Algorithm{AREA}
{Area of Union of Rectangles with Sweep-line and Segment Tree}
{$\mathcal{O}(n \log n)$}
{}{source/pseudocode/area_union_rectangles.txt}

\Algorithm{Manhattan MST}
{Manhattan Minimum Spanning Tree via 8-octant sweep}
{$\mathcal{O}(n \log n)$}
{}{source/pseudocode/manhattan_mst.txt}

\Algorithm{Planar Dual Graph}
{Planar embedding to dual graph: extracts faces (including outer face) and dual edges via angular sorting on half-edges (darts).}
{$\mathcal{O}(E \log E)$}
{}{source/pseudocode/planar_dual.txt}


\section{Tiện ích và kỹ thuật}

\subsection{Template}
\inputminted[]{cpp}{source/template/template.cpp}

\subsection{Debugging}
\inputminted[]{cpp}{source/template/debug.cpp}

\subsection{Script biên dịch và chạy}
Cần \texttt{g++} trong PATH (Windows: MinGW-w64). Truyền tên file nguồn; dừng nếu biên dịch lỗi.
\inputminted[]{}{source/pseudocode/compile_scripts.txt}

\subsection{Phân tích biểu thức (Expression Parsing)}
\inputminted[]{cpp}{source/cpp/utility/expression.cpp}

\subsection{Lập lịch Johnson}
Lập lịch trên hai máy.
\inputminted[]{cpp}{source/cpp/utility/johnson.cpp}

\subsection{Fracturing Search ($k$-best)}
Oracle tìm nghiệm tốt nhất trong miền; heap dùng cùng tiêu chí xếp hạng.
\inputminted[]{}{source/pseudocode/fracturing_search.txt}
\noindent\textbf{Principle}: thứ tự quyết định cố định; child khác $S$ lần đầu tại $i$, prefix giữ các quyết định trước đó giống $S$. Các miền con rời nhau và phủ phần còn lại; không gộp nghiệm theo cost. Chi phí phụ thuộc số miền và oracle.


\subsection{Monte Carlo và Las Vegas}
\label{subsec:randomized}
\begin{itemize}
    \item \textcolor{blue}{\textbf{Monte Carlo}} (\textit{Thời gian cố định, xác suất lỗi $\le\epsilon$; lặp $k$ lần để giảm lỗi xuống $(1-p)^k$}):
    \begin{itemize}
        \item \textbf{Freivalds ($AB\stackrel{?}{=}C$)}: Kiểm tra $P=A(Bv)-Cv$ với $v\in\{0,1\}^N$ ngẫu nhiên trong $O(N^2)$. Nếu $AB\ne C$, $\Pr(P=0)\le1/2$. Lặp $k$ lần $\implies$ lỗi $\le2^{-k}$.
        \item \textbf{Identity Testing / PIT (Schwartz-Zippel)}: Đa thức khác 0 bậc $d$, $P(x_1..x_n)$, tại $r\in S^n$ ngẫu nhiên có $\Pr(P(r)=0)\le d/|S|$. General matching: kiểm tra $\det(\text{Tutte Matrix})\ne0$ trong $O(V^3)$.
        \item \textbf{Fingerprinting / Hashing}: Ánh xạ $v\mapsto r_v\in[0,2^{64}-1]$. Hash tập bằng $\bigoplus r_v$, multiset bằng $\sum cnt(v)r_v\pmod{2^{64}}$, dãy bằng rolling hash 61-bit. Xác suất va chạm $\le2^{-64}$.
        \item \textbf{Color Coding ($k\le8$)}: Tìm đồ thị con/đường đi $k$ đỉnh bằng tô ngẫu nhiên $k$ màu. Xác suất mục tiêu đủ màu là $k!/k^k\ge e^{-k}$. DP trên màu mất $O(2^k(V+E))$. Lặp $O(e^k\ln(1/\epsilon))\approx300$ lần để lỗi $<10^{-6}$.
    \end{itemize}
    \item \textcolor{blue}{\textbf{Las Vegas}} (\textit{Luôn đúng, thời gian ngẫu nhiên, kỳ vọng $1/p$ lần thử}):
    \begin{itemize}
        \item \textbf{Tìm và kiểm chứng witness}: Nếu nghiệm có mật độ $\ge p$, chọn ứng viên ngẫu nhiên rồi kiểm tra trong $O(T)$ $\implies$ thời gian kỳ vọng $O(T/p)$. Ví dụ: chọn phần tử ngẫu nhiên để tìm majority $>n/2$ trong kỳ vọng $O(n)$ ($p>1/2\implies\le2$ lần thử); chọn cặp ngẫu nhiên tìm ước chung.
        \item \textbf{Cấu trúc và chọn ngẫu nhiên}: Pivot QuickSelect (kỳ vọng $O(n)$), priority Treap ngẫu nhiên (chiều cao kỳ vọng $O(\log n)$), Pollard Rho (kỳ vọng $O(n^{1/4})$).
    \end{itemize}
\end{itemize}

\subsection{Lịch (Calendar)}
\inputminted[]{cpp}{source/cpp/utility/calendar.cpp}

\subsection{Josephus}
Đánh số từ 0, mỗi lần loại người thứ $k$: $f(1)=0$, $f(n)=(f(n-1)+k)\bmod n$. Duyệt kích thước $2,\ldots,n$: $O(n)$ thời gian, $O(1)$ bộ nhớ. Với $k=2$: $f(n)=2(n-2^{\lfloor\log_2 n\rfloor})$.

\subsection{Kỹ thuật interactive}
\inputminted[]{cpp}{source/cpp/utility/interactive.cpp}
\begin{itemize}
    \setlength{\itemsep}{0pt}
    \setlength{\parsep}{0pt}
    \item \textbf{Kiểm tra ba phía}: So sánh hai nhóm $A,B$ cùng kích thước; khi có một mục tiêu và baseline bằng nhau đã biết, trả lời xác định mục tiêu ở $A$, $B$ hoặc nhóm còn lại $C$.
    \item \textbf{Gộp truy vấn}: Hỏi $\gcd(X,\prod p_i^{e_i})$ để khôi phục số mũ các nguyên tố đã chọn; tránh tràn tích. Suy ra nhóm cuối khi biết tổng.
    \item \textbf{Chuẩn hóa}: Với $1\le x\le10^d$, tổng chữ số $S(x(10^d-1))=9d$; đưa trạng thái về $9d$, rồi cộng mục tiêu$-9d$.
    \item \textbf{Top hai theo chi phí}: Xáo trộn phần tử còn lại; so với phần tử nhì trước phần tử nhất để ít truy vấn cực đại thật.
    \item \textbf{Đỉnh cực đại/lan truyền}: Cache giá trị đã hỏi, giữ witness cho đỉnh còn khả thi. Dò mìn: số mìn còn lại $=0$ thì các ô chưa biết an toàn; bằng số ô chưa biết thì tất cả là mìn. Tiếp tục trừ các phương trình chồng nhau.
\end{itemize}

\subsection{Giao tiếp hai lần chạy (Run-Twice)}
\textbf{Nhận diện}: Pha 1 gửi thông điệp giới hạn; Pha 2 dùng nó cùng dữ liệu riêng để suy ra đáp án. Phân biệt $S$ khả năng cần ít nhất $\lceil\log_2 S\rceil$ bit. Kiểm tra kênh có giữ thứ tự, mất dữ liệu hay đổi nhãn không.
\begin{itemize}
    \setlength{\itemsep}{0pt}
    \setlength{\parsep}{0pt}
    \item \textbf{Đóng gói}: $0\le x_i<b_i$, gửi $X=\sum_i x_i\prod_{j<i}b_j$; giải mã lần lượt $x_i=X\bmod b_i$, $X\gets\lfloor X/b_i\rfloor$. Ví dụ $(2,1,3)$, cơ số $(3,2,4)$ cho $X=23$. Cần $\prod b_i-1$ vừa miền truyền; \texttt{long long} chứa tối đa 39 chữ số hệ 3.
    \item \textbf{Mất thứ tự}: Gửi $v_i=iB+x_i$ với $0\le x_i<B$, $0\le i<k$; lấy lại $i=\lfloor v_i/B\rfloor$, $x_i=v_i\bmod B$. Ví dụ $7,12,24$ với $B=10$ mã hóa $(7,2,4)$. Cần miền truyền chứa $kB-1$.
    \item \textbf{Mất số, checksum nguyên vẹn}: Mất một số: XOR toàn bộ trước khi gửi, XOR tiếp phần còn lại để khôi phục. Mất hai số: gửi $A=\sum x_i$, $B=\sum x_i^2$; trừ phần còn lại được $s=u+v$, $q=u^2+v^2$. Lấy $d=\sqrt{2q-s^2}=|u-v|$, suy ra $(s\pm d)/2$, chưa biết thứ tự. Dùng số nguyên chính xác, \texttt{\_\_int128} nếu cần; không thay tổng bằng tổng modulo.
    \item \textbf{Hoán vị}: $k$ phần tử phân biệt truyền $k!$ trạng thái, $0\le t<k!$. Hai pha biết cùng tập phần tử, sort làm chuẩn. Còn $m$ phần tử: chọn hạng $\lfloor t/(m-1)!\rfloor$ từ 0, xóa nó, cập nhật $t\gets t\bmod(m-1)!$. Giải mã cộng $c_m(m-1)!$, $c_m$ là số phần tử còn lại nhỏ hơn phần tử đang đọc. Ví dụ $a<b<c$: $bca$ truyền $1\cdot2!+1\cdot1!=3$.
    \item \textbf{Lật một ô báo đích}: Bàn $2^r$ bit, chỉ số $0,\ldots,2^r-1$. Đặt $P=\bigoplus_{b_i=1}i$, lật ô $P\oplus T$; Pha 2 XOR chỉ số bit 1 được $T$. Ví dụ $P=5,T=3$: lật ô 6. Nếu $P=T$, lật ô 0. Cần được tự chọn ô lật và giữ nguyên chỉ số ở hai pha.
\end{itemize}

\subsection{Random, Bit Tricks, Bitset}
\inputminted[]{cpp}{source/cpp/utility/bit_tricks.cpp}

\subsection{Biến đổi bài toán (Transformations)}
\noindent\textbf{Khi tìm hướng tiếp cận}
\begin{itemize}
    \setlength{\itemsep}{0pt}
    \setlength{\parsep}{0pt}
    \setlength{\topsep}{0pt}
    \item Xuất phát từ các đại lượng trong đề.
    \item Thay đổi mức độ cụ thể của suy nghĩ hiện tại (từ tổng quát đến chi tiết, từ trừu tượng đến cụ thể).
    \item \textbf{Xuất phát từ đích, suy luận ngược}: Muốn tính/chứng minh kết quả thì cần biết gì trước? Truy ngược đến khi nối được với dữ kiện trong đề.
    \item \textbf{Xuất phát từ vét cạn}: Viết lời giải hiển nhiên, tìm tính toán lặp lại hoặc dư thừa để tối ưu.
    \item \textbf{Thay đổi cách biểu diễn}: Chuyển sang đồ thị, hình học, đại số hoặc xét phần bù để lộ cấu trúc dễ xử lý hơn.
    \item \textbf{Khai thác ràng buộc đặc biệt}: Mỗi giả thiết có tác dụng gì, vì sao đề cần nó, bỏ đi thì điều gì thay đổi?
    \item \textbf{Liên hệ bài toán đã biết}: Tìm cấu trúc quen thuộc và xác định điểm khác biệt cần xử lý.
    \setlength{\parsep}{0pt}
    \item \textcolor{red}{\textbf{Khi bí}}: Đọc lại đề! Xét trường hợp nhỏ ($N\le6$) để tìm quy luật/bất biến; suy ra độ phức tạp ($N\le20\to2^N,500\to N^3,5000\to N^2,2\cdot10^5\to N\log N$).

    \item \textbf{Tối ưu $\to$ quyết định (Binary Search)}: Nếu điều kiện $P(X)$ đơn điệu, binary search $X$ qua kiểm tra khả thi.

    \item \textbf{Hoán vị toàn cục $\to$ đổi chỗ cục bộ}: Lập luận trao đổi greedy: đổi $(i,j)\to(j,i)$ cải thiện mục tiêu $\iff f(i,j)<f(j,i)$.

    \item \textbf{Ràng buộc $\to$ Graph / Flow / Cut}: 2-SAT, difference constraints ($x_j-x_i\le w$), Project Selection (min-cut), Hall, Dilworth (DAG path cover $=|V|-\text{MaxMatching}$).

    \item \textbf{Đại số $\to$ hình học/tính lồi}: Đường bao trên/dưới, Convex Hull, Slope Trick, Minkowski sum, CHT.

    \item \textbf{Trung bình có trọng số $\to$ Convex Hull}: Xem $\lambda_i$ là tỉ lệ trộn ($\lambda_i\ge0$, $\sum_i\lambda_i=1$). Khi đó $x=\sum_i\lambda_ix_i$ nằm giữa các đầu vào: hai điểm cho đoạn $AB$, ba điểm cho miền tam giác, nhiều điểm cho convex hull. Mọi điểm trong convex hull đều biểu diễn được như vậy. Trong 1D: $x\in[\min x_i,\max x_i]$. Chuẩn hóa trọng số dương: $\lambda_i=w_i/\sum_jw_j$.

    \item \textbf{Tỉ số tổng có trọng số $\to$ trung bình tỉ số}: Với $w_i\ge0$, $b_i>0$, $\frac{\sum_iw_ia_i}{\sum_iw_ib_i}$ là trung bình có trọng số của $a_i/b_i$, trọng số mới $w_ib_i$. Vì vậy nằm trong $[\min_i(a_i/b_i),\max_i(a_i/b_i)]$. Không trọng số: $w_i=1$.

    \item \textbf{Trộn trước $\to$ Jensen}: Với $f$ lồi (hình cái bát), tính sau khi trộn rẻ hơn: $f(\sum_i\lambda_ix_i)\le\sum_i\lambda_if(x_i)$. Đọc là $f(\text{trung bình})\le\text{trung bình của }f$.

    \item \textbf{Đoạn toàn cục $\to$ Divide \& Conquer}: Chia $[l,r]\to[l,m]$ và $[m+1,r]$, cộng đóng góp chéo (CDQ D\&C, D\&C DP).

    \item \textbf{Xóa offline $\to$ thêm (đảo thời gian)}: Đảo thứ tự truy vấn $Q\to1$ để đổi xóa thành thêm DSU. Game DP: đi ngược từ trạng thái kết thúc.

    \item \textbf{Chia miền tìm kiếm (Meet-in-Middle)}: Chia $N=N/2+N/2$, giảm $2^N\to2^{N/2}$ bằng hash table/two pointers.

    \item \textbf{Đúng $k$ nhóm $\to$ hệ số phạt Lagrange}: Tối ưu $cost+\lambda m$, $m$ là số nhóm. Binary search $\lambda$ khi $m$ đơn điệu; cần tính lồi và xử lý hòa để khôi phục đúng $k$. Bài \emph{đếm} đúng $k$ tập dùng Binomial Inversion/PIE từ số cách thuộc ít nhất $k$ tập.

    \item \textbf{Tổng toàn cục $\to$ đóng góp phần tử}: Đổi góc nhìn: cộng đóng góp mỗi phần tử/cạnh/cặp $(u,v)$ trên mọi trạng thái hợp lệ (Fubini).

    \item \textbf{Chia theo ngưỡng ($\sqrt N$ Heavy/Light)}: Nhóm $\le B$: duyệt trực tiếp; $>B$: tối đa $N/B$ nhóm, cache riêng. Cân bằng chi phí theo $B$; khác kỹ thuật Mo.

    \item \textbf{Nâng cấp cuối $\to$ chặn trạng thái}: Bỏ nâng cấp cuối, lập lại kế hoạch còn lại; nếu chi phí tăng $\le$ giá nâng cấp thì bỏ được. Với $L\le N$, $k>\lceil\sqrt N\rceil$: $\lceil L/(k-1)\rceil-\lceil L/k\rceil\le1$. Giá nâng cấp $\ge1$, mỗi điểm dừng giá 1 $\implies$ chặn tài nguyên ở $\lceil\sqrt N\rceil$.

    \item \textbf{Gom nhóm thương (Quotient Grouping)}: $\lfloor N/k\rfloor$ có $O(\sqrt N)$ giá trị khác nhau. Với $1\le l\le N$, nhóm cùng thương kết thúc ở $r=\lfloor N/\lfloor N/l\rfloor\rfloor$; tiếp theo $l=r+1$.

    \item \textbf{Đếm chính xác $\to$ bù/PIE}: Đếm trạng thái hợp lệ qua $\text{tổng}-\text{không hợp lệ}$, Stars \& Bars, PIE.

    \item \textbf{Impartial Game $\to$ Nim Sum}: Chuyển game trên DAG thành XOR các giá trị Sprague-Grundy ($SG(u)=\text{mex}\{SG(v)\}$).

    \item \textbf{Undirected Vertex Geography $\to$ Matching}: Quân bắt đầu tại $s$; hai người luân phiên đi tới hàng xóm chưa thăm, xóa đỉnh vừa rời. Người đầu thắng iff $s$ thuộc mọi maximum matching của $G$, tương đương $\nu(G)>\nu(G-s)$.

    \item \textbf{Tổng quát $\to$ trường hợp đặc biệt}: Giải giới hạn dễ trước ($A_i\in\{0,1\}$, cây, $N\le6$) để tìm quy luật rồi tổng quát hóa.

    \item \textbf{Vét cạn $\to$ Branch \& Bound}: Cắt nhánh DFS bằng cận hợp lệ.

    \item \textbf{Manhattan $\leftrightarrow$ Chebyshev (xoay $45^\circ$)}: $(x,y)\mapsto(x+y,x-y)\implies|x_1-x_2|+|y_1-y_2|=\max(|x'_1-x'_2|,|y'_1-y'_2|)$. Tách hai chiều thành hai trục 1D độc lập.

    \item \textbf{Phân rã max/min và trị tuyệt đối}: $\max(A,B)=\frac{A+B+|A-B|}{2}$, $\min(A,B)=\frac{A+B-|A-B|}{2}$; $|x-a|=\max(x-a,a-x)$. Với một khoảng cách, xét bốn cặp dấu $(\pm x\pm y)$. Để tối thiểu $\sum_i(|x-a_i|+|y-b_i|)$, tách hai tổng 1D rồi chọn median của $a_i$ và $b_i$.

    \item \textbf{Đếm theo ngưỡng/lớp (Fubini)}: $\sum_{i=1}^nA_i=\sum_{v=1}^{\max A}\text{count}(A_i\ge v)$ (cộng giá trị $\longleftrightarrow$ đếm phần tử $\ge v$).

    \item \textbf{Bình phương $\to$ tổng số lẻ}: $x^2=1+3+5+\cdots+(2x-1)=\sum_{i=1}^x(2i-1)$.

    \item \textbf{Dãy chữ số 1 $\to$ tổng cấp số nhân}: $\underbrace{11\ldots1}_{i\text{ chữ số}}=1+10+\cdots+10^{i-1}=\frac{10^i-1}{9}$.

    \item \textbf{Dãy tăng/giảm $\to$ biến không âm}: Bỏ ràng buộc thứ tự bằng chênh lệch kề. Nếu $A_1\le\cdots\le A_n$, đặt $D_i=A_i-A_{i-1}\ge0$; nếu $A_1\ge\cdots\ge A_n$, đặt $D_i=A_{i-1}-A_i\ge0$ ($i\ge2$). Thứ tự nghiêm ngặt cần $D_i\ge1$. Khôi phục $A_i=A_1+\sum_{j=2}^iD_j$ (tăng), hoặc $A_i=A_1-\sum_{j=2}^iD_j$ (giảm).

    \item \textbf{Range Add $\to$ mảng hiệu/Prefix Sum}: $A_i=\sum_{j=1}^iD_j$ ($D_i=A_i-A_{i-1}$); cộng đoạn $[l,r]\to$ hai cập nhật điểm trên $D$. Tổng tĩnh $\to S_r-S_{l-1}$.

    \item \textbf{Cửa sổ cố định $\to$ lớp modulo}: XOR hai phép flip cửa sổ dài $x$ bắt đầu tại $i,i+1$: phần giao triệt tiêu, chỉ đảo $i,i+x$ $\to$ cùng lớp modulo $x$. Mỗi cửa sổ chứa đúng một vị trí mỗi lớp; đặt $P_r=\bigoplus_{i\equiv r\pmod x}s_i$. Flip đảo mọi $P_r$, nên $P_r\oplus P_0$ ($1\le r<x$) bất biến và đủ phân loại trạng thái; từ toàn 0 đạt được $s$ iff mọi $P_r$ bằng nhau. Cue: triệt tiêu hai cửa sổ kề $\to$ hai đầu cách $x$ $\to$ modulo. Cần phép triệt tiêu hợp lệ (XOR, hoặc cộng với phép nghịch đảo); không suy rộng sang gán/max.

    \item \textbf{Thu nhỏ điều kiện: toàn cục $\to$ cục bộ}: Tìm điều kiện trên các phần tử kề/cửa sổ nhỏ tương đương điều kiện cả dãy. Sau sắp xếp, $\min_{i\ne j}|A_i-A_j|=\min_i(A_{i+1}-A_i)$; $A_j-A_i\ge(j-i)D\iff A'_{i+1}\ge A'_i$ ($A'_i=A_i-iD$). Cấp số cộng $\iff$ mọi bộ ba liên tiếp thỏa $A_i+A_{i+2}=2A_{i+1}$; áp dụng ngang/dọc biến hình chữ nhật cấp số cộng thành vùng cửa sổ $3\times3$ hợp lệ (LFAC2 2011A). Đoạn dài $\ge2$ có majority ($>50\%$) luôn chứa witness dài 2 ($x,x$) hoặc 3 ($x,y,x$).

    \item \textbf{Tích $\to$ tổng}: $ab=\binom{a+b}{2}-\binom a2-\binom b2$; $\prod A_i\to\sum\log A_i$ (hoặc log xác suất).

    \item \textbf{Đẳng thức bitwise}: $A+B=(A\oplus B)+2(A\text{ AND }B)=(A\text{ OR }B)+(A\text{ AND }B)$; $A\oplus B=(A\text{ OR }B)-(A\text{ AND }B)$. Lật bit hai số kề: $\sum_{i=1}^n\text{popcount}((i-1)\oplus i)=2n-\text{popcount}(n)$.

    \item \textbf{Kỹ thuật ngẫu nhiên (Monte Carlo / Las Vegas)}: Xem \S\ref{subsec:randomized}. Nghiệm mật độ $p$: Monte Carlo lặp $k$ lần, lỗi $\le(1-p)^k$; Las Vegas kiểm chứng trong kỳ vọng $O(T/p)$. Fingerprinting ($\bigoplus r_v,\sum r_v$) kiểm tra tập/multiset; Color Coding tách đồ thị con; Schwartz-Zippel kiểm tra tồn tại đại số (Tutte).

    \item \textbf{Dirichlet/Pigeonhole (4-Sum $a_x+a_y=a_z+a_w$)}: Các giá trị nguyên dương phân biệt: Có tối đa $2\max A$ giá trị tổng. Nếu $\binom N2>2\max A$, pigeonhole bảo đảm hai cặp trùng tổng. Hash tổng $\to$ cặp chỉ số trong $O(N^2)$; tính phân biệt của các giá trị bảo đảm bốn chỉ số khác nhau.

    \item \textbf{Boyer-Moore Majority Vote}: Tìm phần tử majority ($>n/2$ lần) trong $O(n)$ thời gian, $O(1)$ bộ nhớ bằng một bộ đếm.

    \item \textbf{Không bao nhau $\to$ Antichain (Sperner)}: Chọn các tập con/bitmask sao cho không tập nào chứa tập khác: tối đa $\binom n{\lfloor n/2\rfloor}$. \textbf{Construction}: lấy tất cả tập có đúng $\lfloor n/2\rfloor$ phần tử (cùng popcount nên không bao nhau). Đạt cận khi được chọn từ toàn bộ tập con; có ràng buộc thêm phải kiểm tra lại. Ví dụ: $N$ là tích $n$ prime khác nhau, chọn các ước không chia nhau $\to$ lấy tích của đúng $\lfloor n/2\rfloor$ prime.
\end{itemize}

\noindent\textbf{Invariant: nhận diện và kiểm tra}
\begin{itemize}
    \setlength{\itemsep}{0pt}
    \setlength{\parsep}{0pt}
    \item \textbf{Cách tìm}: Đổi biểu diễn để thao tác thành swap/triệt tiêu. Khác bất biến $\to$ không thể; giống nhau cần chứng minh đủ bằng construction/dạng chuẩn.
    \item \textbf{Tổng / tô màu / modulo}: Chuyển lượng $\to\sum a_i$; đối xứng $\to$ tổng xen kẽ/theo màu. Tìm $w$ sao cho $w\cdot\Delta a=0$ (hoặc $0\pmod m$) với mọi thao tác. XOR cùng lượng hai lần triệt tiêu.
    \item \textbf{Hiệu kề / prefix $\to$ swap}: $a_i'=a_{i-1}+a_{i+1}-a_i$ swap hai hiệu kề. XOR $a_i$ vào hai hàng xóm swap hai prefix XOR $P_{i-1},P_i$. Mọi swap kề được phép $\to$ so multiset; giữ biên $a_1$ (hiệu), $P_n$ (prefix XOR).
    \item \textbf{Hoán đổi $\to$ nhóm vị trí / cặp đối xứng}: Swap từng cặp vị trí: cạnh là swap được phép, độc lập trạng thái $\to$ so multiset mỗi component. Swap cách $k\to$ modulo $k$; nhiều bước nhảy thử gcd, kiểm tra biên. Swap nguyên khối prefix/suffix cùng độ dài $l$, $2l\le n$, giữ thứ tự trong khối $\to$ các cặp đối xứng $(a_i,a_{n+1-i})$ chỉ đổi chỗ/đảo hướng: so multiset các cặp $(\min,\max)$, giữ phần tử giữa nếu $n$ lẻ. Điều kiện đủ khi mọi $1\le l\le\lfloor n/2\rfloor$ được phép (đảo thứ tự $l$ cặp đầu và đảo hướng từng cặp); chỉ cho vài $l$ thì chưa chắc đủ.
    \item \textbf{Parity}: Tìm lượng có tính chẵn/lẻ không đổi: tổng, số bit 1, số vật theo màu, số nghịch thế. Kiểm tra mỗi thao tác làm lượng đó đổi một số chẵn. Với hoán vị phân biệt: parity $=(n-\#\text{cycles})\bmod2$; $k$-cycle tương đương $k-1$ swap, nên $k$ lẻ giữ parity. Mọi 3-cycle sinh đúng hoán vị chẵn; phần tử trùng cho phép đổi parity gán nhãn đích.
    \item \textbf{Cùng lớp dư $\to$ GCD các hiệu}: Đặt $d_i=a_{i+1}-a_i$, $g=\gcd_{i<n}|d_i|$. Khi đó $g=\gcd_{i>1}|a_i-a_1|=\gcd_{i<j}|a_i-a_j|$ (hiệu cặp là tổng hiệu kề). Với $g>0$, mọi $a_i$ nằm trên lưới $a_1+g\mathbb Z$; $g$ là modulo lớn nhất để tất cả cùng dư. Tịnh tiến cả mảng hoặc đổi thứ tự không đổi $g$. Cue: khoảng cách giữa giá trị/cùng lớp dư $\to$ xét gcd hiệu, không phải gcd giá trị. Thêm $2u-v$ giữ lưới vì $u,v\equiv a_1\pmod g\implies2u-v\equiv a_1\pmod g$ (CF1477A). $g=0$: mọi giá trị bằng nhau, xét riêng.
\end{itemize}

\noindent\textbf{Monovariant: mỗi bước tiến gần hơn đến đích}
\begin{itemize}
    \setlength{\itemsep}{0pt}
    \setlength{\parsep}{0pt}
    \item \textbf{Cách dùng}: Chọn một số đo $\Phi$ (số lỗi, khoảng cách, độ lệch); chứng minh sau mỗi bước nó giảm. Nếu $\Phi$ là số nguyên không âm và giảm ít nhất 1, quá trình dừng sau tối đa $\Phi$ ban đầu bước. Dừng chưa chắc đạt đích: cần chứng minh chưa đạt đích thì vẫn chọn được bước làm giảm $\Phi$.
    \item \textbf{Sửa thứ tự $\to$ số nghịch thế}: Swap hai phần tử kề $a_i>a_{i+1}$ làm số cặp sai thứ tự giảm đúng 1. Số vị trí sai không chắc giảm khi swap.
    \item \textbf{Di chuyển $\to$ khoảng cách/vị trí}: Mỗi bước đưa vật gần đích hơn $\to$ tổng khoảng cách giảm. Vật chỉ dịch trái $\to$ tổng vị trí của các vật giảm, dù số vật giữ nguyên.
    \item \textbf{San bằng $\to$ tổng bình phương}: Chuyển 1 từ số lớn $u$ sang số nhỏ $v$, $u\ge v+2$: $(u-1)^2+(v+1)^2-u^2-v^2=-2(u-v-1)<0$. Tổng không đổi, nhưng độ lệch giảm.
    \item \textbf{Cố định dần $\to$ số chỗ còn phải sửa}: Mỗi bước kéo dài prefix đúng và không phá nó $\to$ số vị trí ngoài prefix giảm. Nếu số đo chính đôi khi không đổi, thêm số đo phụ: bước nào cũng giảm số chính, hoặc giữ số chính và giảm số phụ (cả hai nguyên không âm).
    \item \textbf{Kiểm tra}: Chỉ ``không tăng'' có thể lặp; số thực giảm như $1,1/2,1/4,\ldots$ cũng có thể vô hạn. Nếu chỉ các bước mình chọn làm $\Phi$ giảm, kết luận chỉ áp dụng cho thuật toán đó.
\end{itemize}


\begin{itemize}
    \setlength{\itemsep}{0pt}
    \setlength{\parsep}{0pt}
    \item \textbf{Chia miền $\to$ giảm một nửa số mũ (LFAC2 2014H, 4 sao)}: Tem RGB ngắn nhất phủ chuỗi có $k$ wildcard; chồng tại wildcard được khác màu. $L\ge n/2$: chỉ cần prefix/suffix tương thích. $L\le n/2$: hai đầu dài $L$ rời nhau, chọn đầu ít wildcard hơn, tối đa $\lfloor k/2\rfloor$; liệt kê $3^{k/2}$ thay vì $3^k$.

    \item \textbf{Bộ ba Pythagoras}: Nguyên thủy: $(a,b,c)=(u^2-v^2,2uv,u^2+v^2)$, $u>v\ge1$, $\gcd(u,v)=1$, khác parity. Sinh mọi bộ bằng lấy bội $k(a,b,c)$; nếu $c\le C$, duyệt $u^2+v^2\le C$, $k\le C/c$.

    \item \textbf{Factorial: nhánh nhỏ và nhánh khó (CRRC K)}: Tìm $a!/(b!c!)=N$, chọn $b\ge c$, đặt $d=a-b$. Nếu $d\ge c$, $N\ge\binom{2c}{c}$ chặn $c$; nhánh $d<c$ phải khai thác chia hết theo prime valuations. Với $N\le10^9$, editorial dùng $d\le20$; cutoff thực nghiệm, chưa chứng minh ở đây.
\end{itemize}


\subsection{Phím tắt VS Code (Linux)}
\textbf{Mặc định}; dấu phẩy: bấm hai tổ hợp lần lượt. \texttt{Ctrl+K, Ctrl+S}: tra/gán phím nếu keymap khác hoặc hệ điều hành chặn.
\begin{itemize}
    \setlength{\itemsep}{0pt}
    \setlength{\parsep}{0pt}
    \item \textbf{Nhiều con trỏ}: \texttt{Ctrl+D} chọn lần xuất hiện kế; \texttt{Ctrl+K, Ctrl+D} bỏ qua lần vừa chọn; \texttt{Ctrl+U} hoàn tác con trỏ; \texttt{Ctrl+Shift+L} chọn mọi lần xuất hiện của vùng chọn; \texttt{Esc} thoát.
    \item \textbf{Đổi tên}: \texttt{F2} đổi tên symbol theo scope (cần hỗ trợ ngôn ngữ); \texttt{Ctrl+H} thay chuỗi trong file.
    \item \textbf{Nhiều dòng}: \texttt{Alt+Click} thêm con trỏ; \texttt{Shift+Alt+I} thêm con trỏ cuối các dòng được chọn (ví dụ thêm dấu phẩy hàng loạt).
    \item \textbf{Dòng}: \texttt{Alt+}$\uparrow/\downarrow$ chuyển dòng/khối; \texttt{Ctrl+Shift+K} xóa dòng. Không chọn gì: \texttt{Ctrl+C/X} copy/cắt dòng.
    \item \textbf{Chèn / format}: \texttt{Ctrl+Enter} dòng mới dưới; \texttt{Ctrl+Shift+Enter} dòng mới trên; \texttt{Ctrl+/} comment dòng; \texttt{Ctrl+Shift+I} format file (cần formatter).
    \item \textbf{Nhảy}: \texttt{Ctrl+G} tới dòng; \texttt{Ctrl+Shift+O} tìm symbol; \texttt{F12} tới định nghĩa; \texttt{Ctrl+Shift+\textbackslash} tới ngoặc khớp; \texttt{F8} lỗi/cảnh báo kế.
    \item \textbf{File / lệnh}: \texttt{Ctrl+P} tìm file; \texttt{Ctrl+Tab} đổi file; \texttt{Ctrl+Shift+T} mở lại tab đóng; \texttt{Ctrl+Shift+P} tìm lệnh.
    \item \textbf{Terminal}: \texttt{Ctrl+\textasciigrave} hiện/ẩn; \texttt{Ctrl+Shift+\textasciigrave} tạo mới; \texttt{Ctrl+Shift+C/V} copy/paste trong terminal.
\end{itemize}
\textit{Nguồn: VS Code, Keyboard shortcuts for Linux. Alt+Click bị chặn: đổi Multi Cursor Modifier.}

\subsection{Checklist trước khi nộp}
\begin{itemize}
    \setlength{\itemsep}{0pt}
    \setlength{\parsep}{0pt}
    \item \textcolor{red}{\textbf{Đã đọc giới hạn và edge case ($N=0,1$, âm, mất liên thông)?}}
    \item Đã reset/xóa trạng thái toàn cục giữa các test case?
    \item Đúng giới hạn mảng, chỉ số từ 0 hay 1?
    \item Có tràn số nguyên (dùng \texttt{long long}/\texttt{\_\_int128})?
    \item Hàm thiếu \texttt{return} (Undefined Behavior)?
    \item Độ chính xác số thực/so sánh epsilon (\texttt{eps = 1e-9})?
    \item Có fast I/O (\texttt{cin.tie(NULL)}, \texttt{'\textbackslash n'} thay \texttt{endl})?
\end{itemize}

\end{multicols}
\end{document}
